\chapter{KL 估计的统一残差框架}
\label{chap:kl-estimation}

KL 的数值、KL 的参数梯度和某段 loss 的自动微分梯度是三个对象。本章先把前两个对象写成有限词表上的积分，再用同一个残差校正定理分析估计误差，最后检查自动微分是否实现了所需的梯度。算法名称放在相应积分与控制函数确定之后；无偏性和方差由这些对象推出。

\section{目标、采样测度与微分规则}
\label{sec:kl-targets}
\subsection{固定条件与唯一符号}

固定 prefix $x$，省略条件变量。词表 $\mathcal V$ 有限，$p=p_\theta$ 是当前 student，$q$ 是冻结参考，$\rho$ 是实际动作采样分布，$A\sim\rho$。$q$ 默认代入 teacher；分析 student--anchor KL 时只代入 $q=\pi_k$。冻结 rollout 则代入 $\rho=\pi_k$。参考分布与采样分布具有不同职责。

假设 $p(v),q(v)>0$。对 $\rho$ 的支持要求由残差决定，并在总定理中给出。所有期望、协方差和偏差先条件于固定 prefix；若使用历史数据拟合控制，也先条件于这些数据。只对 $p$ 求导，$q$ 与支持集在求导时冻结。统一定义
\[
 r_v=\log\frac{p(v)}{q(v)},\qquad
 s_v=\nabla_\theta\log p(v),\qquad
 \sg(\cdot)=\operatorname{stopgrad}(\cdot).
\]
$r_v$ 始终是 KL log-ratio；前两章的冻结 OPD reward 保留完整记号 $r_k(x,v)$，在 $q=\pi_T$ 时它对应负 log-ratio。$s_v$ 始终表示 score function（得分函数）。

\begin{definition}[积分与误差]
对概率测度 $\nu$ 及标量或向量函数 $f$，记
\[
 I_\nu(f)=\sum_{v\in\mathcal V}\nu(v)f_v.
\]
这里 $\nu$ 只标记求和权重，代入时取 $p$ 或 $q$，没有引入额外策略。目标 $I_\nu(f)$ 的估计量 $\widehat I$ 的偏差、协方差迹与均方误差为
\[
 \operatorname{Bias}(\widehat I)=\E_\rho\widehat I-I_\nu(f),\qquad
 \Var_\rho(\widehat I)=\operatorname{tr}\operatorname{Cov}_\rho(\widehat I),
\]
\[
 \operatorname{MSE}(\widehat I)
 =\|\operatorname{Bias}(\widehat I)\|^2+\Var_\rho(\widehat I).
\]
标量情形把范数解释为绝对值。估计数值和估计梯度时，只替换 $f$ 的值域。
\end{definition}

\subsection{四个目标积分}
记反向和正向 KL 为
\[
 K=\KL(p\Vert q),\qquad K^{\mathrm F}=\KL(q\Vert p).
\]
有限求和的乘积法则及归一化恒等式给出
\[
 \sum_vp(v)s_v=\nabla_\theta\sum_vp(v)=0,
\]
\begin{align*}
 g=\nabla_\theta K
 &=\sum_v\{(\nabla p(v))r_v+p(v)\nabla r_v\}\\
 &=\sum_vp(v)(r_v+1)s_v=\sum_vp(v)r_vs_v,\\
 g^{\mathrm F}=\nabla_\theta K^{\mathrm F}
 &=-\sum_vq(v)\nabla r_v=-\sum_vq(v)s_v.
\end{align*}
因此整章只需要下列四次代入：
\begin{center}
\begin{tabular}{llll}
\toprule
目标 & 积分权重 $\nu$ & 被积函数 $f_v$ & 积分值\\
\midrule
反向 KL 数值 & $p$ & $r_v$ & $K=I_p(r)$\\
反向 KL 梯度 & $p$ & $r_vs_v$ & $g=I_p(rs)$\\
正向 KL 数值 & $q$ & $-r_v$ & $K^{\mathrm F}=I_q(-r)$\\
正向 KL 梯度 & $q$ & $-s_v$ & $g^{\mathrm F}=I_q(-s)$\\
\bottomrule
\end{tabular}
\end{center}
积分权重 $\nu$ 决定目标，实际采样分布 $\rho$ 决定怎样估计目标。重要性比率是 $\nu/\rho$，由积分权重与实际采样分布共同决定。

\subsection{数值微分与随机梯度的区别}
对可微函数 $F_\theta(v)$，逐项求导得到
\begin{equation}\label{eq:kl-measure-derivative}
 \nabla_\theta I_p(F_\theta)
 =I_p\bigl(\nabla_\theta F_\theta+F_\theta s\bigr).
\end{equation}
固定样本索引后，自动微分只能得到 $\nabla_\theta F_\theta(A)$；采样权重的微分项 $F_\theta(A)s_A$ 不会由离散抽样自动出现。反向 KL 取 $F_\theta=r$，两项分别为 $s$ 与 $rs$，前者的总体均值恰好为零。第~\ref{sec:kl-autodiff} 节将使用这条恒等式统一审计 K1--K5 与稀疏头尾反传。

本章主线是：先选上表的目标，随后选可精确积分的近似 $a$，用随机样本补回 $f-a$，再选择采样分布和反传 surrogate。确定性截断对应丢弃残差；HETS 对头部取精确近似；ResOPD 与 SC-TA 对尾部采用具有相同积分、不同逐 token 值的近似。

\section{统一构造：精确积分加残差抽样}
\label{sec:kl-residual-framework}

KL 数值、KL 梯度和 baseline 的统计分析可以共用一个积分问题。固定当前 prefix 后，$p$、$q$、$r_v$ 与 $s_v$ 都是词表上的确定函数。选择积分测度 $\nu$ 和被积函数 $f$，便确定了估计目标；选择实际抽样分布 $\rho$ 和可精确积分的近似函数 $a$，便确定了估计方法。以下所有协方差均针对动作抽样，向量的方差仍指协方差迹。

\begin{definition}[可精确积分的近似函数与残差]
\label{def:kl-residual-integral}
令 $\nu\in\{p,q\}$，$f:\mathcal V\to\mathbb R^d$，并定义
\[
 R_v=f_v-a_v,\qquad I_\nu(f)=I_\nu(a)+I_\nu(R).
\]
其中 $a:\mathcal V\to\mathbb R^d$ 的积分 $I_\nu(a)$ 可以精确计算。标量情形取 $d=1$。令 $A\sim\rho$，假设
\[
 \nu(v)\lVert R_v\rVert>0\ \Longrightarrow\ \rho(v)>0.
\]
这个支持条件只要求覆盖非零残差。定义残差修正估计量
\begin{equation}
 \widehat I_{\nu,\rho,a}(f)
 =I_\nu(a)+\frac{\nu(A)}{\rho(A)}R_A.
 \label{eq:kl-master-estimator}
\end{equation}
比率只在实际抽到的动作上计算。$\rho(v)=0$ 且残差贡献为零的动作无需定义该比率。
\end{definition}

该定义中，$\nu$ 决定要计算哪个积分，$\rho$ 决定从哪里取得随机动作，$a$ 决定哪些贡献已经被确定性地计算。三个选择应当分别记录。特别地，参考分布 $q$ 出现在 $r_v$ 内，并不意味着积分测度或抽样测度就是 $q$。

本节条件于当前参数，把这些函数当作词表上的确定值。估计梯度时直接代入向量函数 $f_v=r_vs_v$ 或 $-s_v$；对一个无偏数值估计量再次自动微分，需要另外验证第~\ref{sec:kl-autodiff} 节的规则。

\begin{theorem}[残差修正的均值、协方差与均方误差]
\label{thm:kl-master-moments}
记 $m_R=I_\nu(R)$。在定义~\ref{def:kl-residual-integral} 的支持条件下，式~\eqref{eq:kl-master-estimator} 无偏，并且
\begin{align}
 \operatorname{Cov}(\widehat I_{\nu,\rho,a}(f))
 &=\sum_{v:\rho(v)>0}\frac{\nu(v)^2}{\rho(v)}R_vR_v^\top-m_Rm_R^\top,
 \label{eq:kl-master-cov}\\
 \Var(\widehat I_{\nu,\rho,a}(f))
 &=\sum_{v:\rho(v)>0}\frac{\nu(v)^2}{\rho(v)}\lVert R_v\rVert^2-\lVert m_R\rVert^2.
 \label{eq:kl-master-trace}
\end{align}
其相对于 $I_\nu(f)$ 的均方误差等于式~\eqref{eq:kl-master-trace}。$n$ 个独立同分布动作对应的样本均值，其协方差和均方误差分别为上述结果的 $1/n$。
\end{theorem}

\begin{proof}
先逐项计算随机修正项的期望。支持条件保证被略去的动作对 $I_\nu(R)$ 没有贡献，因此
\[
 \E_\rho\!\left[\frac{\nu(A)}{\rho(A)}R_A\right]
 =\sum_{v:\rho(v)>0}\nu(v)R_v=I_\nu(R).
\]
加上 $I_\nu(a)$，得到 $I_\nu(a+R)=I_\nu(f)$。确定项不影响协方差；令 $Y=\nu(A)R_A/\rho(A)$，有
\[
 \E[YY^\top]
 =\sum_{v:\rho(v)>0}\frac{\nu(v)^2}{\rho(v)}R_vR_v^\top,
 \qquad \E[Y]=m_R.
\]
应用 $\operatorname{Cov}(Y)=\E[YY^\top]-\E[Y]\E[Y]^\top$ 得到矩阵公式。对矩阵取迹，并使用 $\operatorname{tr}(zz^\top)=\lVert z\rVert^2$，得到迹方差。最后，无偏性给出
\[
 \E\lVert\widehat I-I_\nu(f)\rVert^2
 =\operatorname{tr}\operatorname{Cov}(\widehat I).
\]
独立样本之间的交叉协方差为零，因而平均后的协方差除以 $n$。
\end{proof}

\subsection{精确项误差、错误权重与缺失支持}

残差公式同时给出实现偏差的核对方式。令 $\mathcal D=\{v:\rho(v)>0\}$，允许实际使用的权重 $\widetilde w(v)$ 偏离 $\nu(v)/\rho(v)$，并允许确定积分项存在固定误差 $\delta\in\mathbb R^d$。实际输出为
\[
 \widetilde I=I_\nu(a)+\delta+\widetilde w(A)R_A.
\]
这里先条件于所有已冻结的缓存和控制参数；随机缓存的处理在本节末给出。

\begin{theorem}[三种误差来源的统一偏差]
\label{thm:kl-master-bias}
上述输出相对于 $I_\nu(f)$ 的偏差恰为
\begin{equation}
 b_I=\delta
 +\sum_{v\in\mathcal D}\bigl(\rho(v)\widetilde w(v)-\nu(v)\bigr)R_v
 -\sum_{v\notin\mathcal D}\nu(v)R_v.
 \label{eq:kl-master-bias}
\end{equation}
记 $\widetilde m_R=\sum_{v\in\mathcal D}\rho(v)\widetilde w(v)R_v$，则
\begin{align}
 \operatorname{Cov}(\widetilde I)
 &=\sum_{v\in\mathcal D}\rho(v)\widetilde w(v)^2R_vR_v^\top
 -\widetilde m_R\widetilde m_R^\top,\\
 \E\lVert\widetilde I-I_\nu(f)\rVert^2
 &=\operatorname{tr}\operatorname{Cov}(\widetilde I)+\lVert b_I\rVert^2.
 \label{eq:kl-master-mse}
\end{align}
\end{theorem}

\begin{proof}
由 $I_\nu(f)=I_\nu(a)+\sum_v\nu(v)R_v$，直接相减得
\[
 \E[\widetilde I]-I_\nu(f)
 =\delta+\sum_{v\in\mathcal D}\rho(v)\widetilde w(v)R_v-\sum_v\nu(v)R_v.
\]
把最后一个求和分成 $\mathcal D$ 及其补集，就得到式~\eqref{eq:kl-master-bias}。协方差由随机项的一阶、二阶矩计算。均方误差展开为
\[
 \E\lVert(\widetilde I-\E\widetilde I)+b_I\rVert^2
 =\E\lVert\widetilde I-\E\widetilde I\rVert^2+\lVert b_I\rVert^2,
\]
其中交叉项因 $\E[\widetilde I-\E\widetilde I]=0$ 消失。
\end{proof}

精确积分误差、重要性权重误差和缺失支持具有不同来源，却都在同一个残差上累积。若近似函数已经精确覆盖某个动作，即 $R_v=0$，该动作的权重误差或支持缺失不会产生偏差。反之，即使头部计算完全精确，缺失尾部的残差均值仍会留下系统偏差。无偏性的严格条件是式~\eqref{eq:kl-master-bias} 为零；三个误差项也能相互抵消，因此分别为零是便于验证的充分条件。

\subsection{已知积分控制函数的最优线性组合}

现在固定 $\nu$ 与 $\rho$，假设 $\rho$ 覆盖 $\nu$ 的支持。取 $m$ 个向量控制函数 $h_j:\mathcal V\to\mathbb R^d$，每个积分 $I_\nu(h_j)$ 已知。标量系数 $c_j$ 组成 $c\in\mathbb R^m$，并令
\[
 a_{c,v}=\sum_{j=1}^m c_jh_{j,v},\qquad
 w(v)=\frac{\nu(v)}{\rho(v)},\qquad X=w(A)f_A,
\]
\[
 U_j=w(A)h_{j,A}-I_\nu(h_j),\qquad
 X_0=X-I_\nu(f).
\]
于是 $\E U_j=0$，且主估计量可以逐样本写成
\[
 \widehat I_c=X-\sum_jc_jU_j.
\]
这一步把选择近似函数转化为在零均值随机函数空间中投影。

\begin{theorem}[Gram 矩阵、最优控制与精确超额方差]
\label{thm:kl-master-gram}
定义对称半正定矩阵 $\Gamma$、向量 $\zeta$ 和基准方差 $V_0$：
\[
 \Gamma_{ij}=\E_\rho\langle U_i,U_j\rangle,\qquad
 \zeta_j=\E_\rho\langle U_j,X_0\rangle,\qquad
 V_0=\E_\rho\lVert X_0\rVert^2.
\]
令 $\Gamma^\dagger$ 为 Moore--Penrose 伪逆，则迹方差最优的系数集合为
\[
 c^*=\Gamma^\dagger\zeta+\ker\Gamma.
\]
最优值及任意系数相对于最优值的代价分别为
\begin{align}
 V_{\min}&=V_0-\zeta^\top\Gamma^\dagger\zeta,\\
 \Var(\widehat I_c)-V_{\min}
 &=(c-\Gamma^\dagger\zeta)^\top\Gamma(c-\Gamma^\dagger\zeta).
 \label{eq:kl-master-gram-regret}
\end{align}
这里最小化的是协方差迹；一般向量控制下不同时最小化所有方向的方差。
\end{theorem}

\begin{proof}
估计量的中心化形式为 $X_0-\sum_jc_jU_j$，故展开平方得到
\[
 \Var(\widehat I_c)=V_0-2c^\top\zeta+c^\top\Gamma c.
\]
对任意 $z\in\ker\Gamma$，有
\[
 0=z^\top\Gamma z=\E\left\lVert\sum_jz_jU_j\right\rVert^2.
\]
因此 $\sum_jz_jU_j=0$ 几乎处处成立，进而
\[
 z^\top\zeta=\E\left\langle\sum_jz_jU_j,X_0\right\rangle=0.
\]
这证明 $\zeta\perp\ker\Gamma$，即 $\zeta\in\operatorname{range}\Gamma$。令 $c_0=\Gamma^\dagger\zeta$，便有 $\Gamma c_0=\zeta$。对二次式配方得
\[
 V_0-2c^\top\zeta+c^\top\Gamma c
 =V_0-c_0^\top\Gamma c_0+(c-c_0)^\top\Gamma(c-c_0).
\]
由伪逆恒等式，$c_0^\top\Gamma c_0=\zeta^\top\Gamma^\dagger\zeta$。最后一个二次型非负，且恰在 $c-c_0\in\ker\Gamma$ 时为零，得到全部结论。
\end{proof}

单控制的协方差除以方差公式是该定理在 $m=1$ 时的特例；多个相关控制必须使用整个 Gram 矩阵。控制之间的相关项决定一个控制能否补充另一个控制已经消去的波动。增加控制函数使可选空间包含原来的空间，因此最优迹方差只能下降或保持不变；实际使用不准确的系数，其代价由式~\eqref{eq:kl-master-gram-regret} 精确给出。

\subsection{固定残差时的最优抽样分布}

控制函数改变残差，抽样分布改变残差的二阶矩权重。因此控制优化与抽样优化属于同一方差问题的两个坐标。

\begin{theorem}[残差范数决定迹方差最优抽样]
\label{thm:kl-master-proposal}
固定 $a$，令 $L_R=\sum_v\nu(v)\lVert R_v\rVert$。若 $L_R>0$，则主估计量的最优抽样分布为
\[
 \rho^*(v)=\frac{\nu(v)\lVert R_v\rVert}{L_R},
\]
其最小迹方差为
\[
 \Var_{\min}=L_R^2-\lVert I_\nu(R)\rVert^2.
\]
若 $L_R=0$，则 $I_\nu(a)=I_\nu(f)$，直接返回确定积分项便有零方差。
\end{theorem}

\begin{proof}
令 $\mathcal R=\{v:\nu(v)\lVert R_v\rVert>0\}$。对任何满足残差支持条件的 $\rho$，由 Cauchy--Schwarz 不等式，
\begin{align*}
 L_R^2
 &=\left(\sum_{v\in\mathcal R}\frac{\nu(v)\lVert R_v\rVert}{\sqrt{\rho(v)}}\sqrt{\rho(v)}\right)^2\\
 &\le\left(\sum_{v\in\mathcal R}\frac{\nu(v)^2\lVert R_v\rVert^2}{\rho(v)}\right)
       \left(\sum_{v\in\mathcal R}\rho(v)\right)\\
 &\le\sum_{v\in\mathcal R}\frac{\nu(v)^2\lVert R_v\rVert^2}{\rho(v)}.
\end{align*}
取 $\rho=\rho^*$ 时，全部质量放在 $\mathcal R$，且两向量成比例，两个不等式均取等。代入式~\eqref{eq:kl-master-trace} 得到最优值。若 $L_R=0$，则 $\nu(v)R_v=0$ 对所有动作成立，目标已被精确积分项覆盖。
\end{proof}

在标量情形，最优二阶矩由 $\lvert R\rvert$ 决定；在向量情形，由欧氏范数决定。即使取最优抽样，残差方向随动作变化仍会产生方差。零方差要求三角不等式取等，也就是所有非零残差向量沿同一个有向射线。若额外要求 $\rho$ 对整个词表严格为正，而最优解对零残差动作赋零质量，则上述值是可逼近的下确界。

改变 $a$ 会同时改变 $\Gamma$、$\zeta$ 和 $\rho^*$；改变 $\rho$ 也会改变控制系数的最优解。因而在某个采样分布上拟合出的最优控制系数，不能直接声明在另一个采样分布上仍然最优。

这些最优值是总体矩已知时的 oracle 结论。稀疏 teacher payload 不提供所有未知尾部的 $r_v$，因此一般不能直接构造 $\rho^*$ 或精确计算最优 Gram 系数。它们用于给出可达到的下界、判断控制特征是否有用，以及设计独立样本上的近似拟合；可计算算法必须另外说明这些量如何取得。

\subsection{分区后统一分离事件内部与事件之间的方差}

\begin{theorem}[on-policy 残差的分区协方差]
\label{thm:kl-master-partition}
令 $\rho=\nu$，$\{C_\ell\}_{\ell=1}^L$ 为词表分区，忽略零质量单元。记
\[
 \omega_\ell=\nu(C_\ell),\quad
 \mu_\ell=\E_\nu[R_A\mid A\in C_\ell],\quad
 \Sigma_\ell=\operatorname{Cov}_\nu(R_A\mid A\in C_\ell),\quad
 \bar\mu=\sum_\ell\omega_\ell\mu_\ell.
\]
则
\begin{equation}
 \operatorname{Cov}(\widehat I_{\nu,\nu,a}(f))
 =\sum_\ell\omega_\ell\Sigma_\ell
  +\sum_\ell\omega_\ell(\mu_\ell-\bar\mu)(\mu_\ell-\bar\mu)^\top.
 \label{eq:kl-master-partition}
\end{equation}
\end{theorem}

\begin{proof}
因为重要性比率为一，估计量的协方差就是 $R_A$ 的协方差。条件于 $A\in C_\ell$，写成
\[
 R_A-\bar\mu=(R_A-\mu_\ell)+(\mu_\ell-\bar\mu).
\]
外积展开后，第一项的条件期望为 $\Sigma_\ell$，第二项为确定外积。两个交叉项均含 $\E[R_A-\mu_\ell\mid A\in C_\ell]=0$，故为零。最后按单元概率 $\omega_\ell$ 求和即得结论。
\end{proof}

该公式使头尾控制的作用可以逐项识别。在 $\nu=\rho=p$ 的例子中，取头部 $S$ 与尾部 $S^c$，局部记 $P=p(S)$、$\tau_p=1-P$、$J_A=\mathbf1\{A\notin S\}$。令 $a_v=f_v$ 在头部成立，且 $a_v=b$ 在尾部成立。记尾部条件均值为 $\mu_f$、条件协方差为 $\Sigma_f$，则
\[
 I_p(a)=\sum_{v\in S}p(v)f_v+\tau_pb,\qquad
 R_A=J_A(f_A-b),
\]
\begin{equation}
 \operatorname{Cov}(\widehat I)
 =\tau_p\Sigma_f+P\tau_p(\mu_f-b)(\mu_f-b)^\top.
 \label{eq:kl-master-head-tail}
\end{equation}
减去常向量 $b$ 不改变尾部内部协方差，只改变头尾事件之间的条件均值差。取 $b=\mu_f$ 消去事件间项；要进一步降低事件内部项，需要尾部随动作变化的控制函数。数值目标与梯度目标只需分别代入 $f_v=r_v$ 与 $f_v=r_vs_v$，全部计算便沿用同一分区公式。

\subsection{残差均值、残差波动与随机控制的条件化}

若仅返回确定近似 $I_\nu(a)$，其偏差为 $-I_\nu(R)$，方差为零。若补上正确的残差抽样，偏差变为零，留下的是加权残差的中心二阶矩。由此，近似误差的均值与抽样误差的波动必须分别分析：残差均值为零不保证方差小，残差几乎为常向量也不保证确定近似无偏。

控制函数、系数和采样分布可以由历史数据学习。令 $\mathcal F$ 包含拟合数据和冻结缓存；若条件于 $\mathcal F$，这些对象已固定，且新动作满足 $A\mid\mathcal F\sim\rho_{\mathcal F}$，前述定理均可逐条件应用。若目标 $I_\nu(f)$ 本身固定且每个条件估计量无偏，则
\[
 \E[\widehat I\mid\mathcal F]=I_\nu(f),\qquad
 \operatorname{Cov}(\widehat I)=\E[\operatorname{Cov}(\widehat I\mid\mathcal F)].
\]
若条件偏差为随机向量 $b_{\mathcal F}$，全协方差公式则给出
\[
 \operatorname{Cov}(\widehat I)
 =\E[\operatorname{Cov}(\widehat I\mid\mathcal F)]+\operatorname{Cov}(b_{\mathcal F}),
 \qquad \operatorname{Bias}(\widehat I)=\E[b_{\mathcal F}].
\]
同一动作同时参与控制拟合与残差评估时，条件抽样关系需要重新证明。例如 $\E U=0$ 一般不能推出 $\E[c(A)U]=0$。独立校准集或交叉拟合（cross-fitting）可以使每个评估动作独立于用于构造其控制函数的训练数据，从而恢复逐条件无偏性；交叉拟合各折之间仍能相关，总方差不能直接按独立样本数相除。

\section{同一分区上的 KL 粗化与梯度残差}
\label{sec:kl-partition-geometry}
\label{sec:kl-partition}

头部精确求和、尾部合并和尾部残差抽样，都可以从同一个词表分区推导。分区同时规定可精确计算的粗粒度目标和需要抽样的细粒度残差。下面始终固定 prefix、当前参数和参考分布；只在明确写出梯度时对 student 参数求导。

\begin{definition}[词表分区与桶内分布]
令 $\mathcal C$ 是有限词表 $\mathcal V$ 的一个分区，各个单元 $C\in\mathcal C$ 非空且两两不交，并满足 $\bigcup_{C\in\mathcal C}C=\mathcal V$。分区在求导时固定。定义
\[
 P_C=\sum_{v\in C}p(v),\qquad
 Q_C=\sum_{v\in C}q(v),\qquad
 \ell_C=\log\frac{P_C}{Q_C},
\]
以及桶内条件分布和平均 score
\[
 p_C(v)=\frac{p(v)}{P_C},\qquad
 q_C(v)=\frac{q(v)}{Q_C},\qquad
 \bar s_C=\sum_{v\in C}p_C(v)s_v=\nabla\log P_C.
\]
用 $C(v)$ 表示包含 $v$ 的唯一单元。粗粒度反向 KL 与桶内形状 KL 分别定义为
\[
 K_{\mathcal C}=\sum_{C\in\mathcal C}P_C\ell_C,
 \qquad D_C=\KL(p_C\Vert q_C).
\]
这里 $K_{\mathcal C}$ 是桶标签分布之间的 KL；$D_C$ 描述保留桶标签以后仍未解释的条件分布差异。
\end{definition}

\begin{theorem}[同一分区上的双向 KL 链式分解]
\label{thm:kl-partition-chain}
在本章正支持假设下，反向 KL 和正向 KL 分别满足
\begin{align}
 K&=K_{\mathcal C}+\sum_{C\in\mathcal C}P_CD_C,
 \label{eq:kl-partition-chain-reverse}\\
 K^{\mathrm F}&=K^{\mathrm F}_{\mathcal C}
       +\sum_{C\in\mathcal C}Q_CD^{\mathrm F}_C,
 \label{eq:kl-partition-chain-forward}
\end{align}
其中
\[
 K^{\mathrm F}_{\mathcal C}=-\sum_CQ_C\ell_C,
 \qquad D_C^{\mathrm F}=\KL(q_C\Vert p_C).
\]
因此两种方向的粗粒度 KL 都是对应完整 KL 的下界。等号成立当且仅当每个单元内的条件分布相同。
\end{theorem}

\begin{proof}
对任意 $v\in C$，完整 log-ratio 可以逐点写为
\[
 r_v=\log\frac{p(v)}{q(v)}
     =\ell_C+\log\frac{p_C(v)}{q_C(v)}.
\]
先乘以 $p(v)=P_Cp_C(v)$，在单元内求和，再对单元求和，得到
\[
 \sum_vp(v)r_v
 =\sum_CP_C\ell_C
  +\sum_CP_C\sum_{v\in C}p_C(v)
                        \log\frac{p_C(v)}{q_C(v)}.
\]
这就是反向分解。将上面的逐点恒等式取负并乘以 $q(v)=Q_Cq_C(v)$，按相同顺序求和，得到正向分解。所有质量严格为正，且每个条件 KL 非负；条件 KL 等于零恰好对应桶内条件分布相同，故得到下界及等号条件。
\end{proof}

\begin{proposition}[分区细化的标量单调性]
\label{prop:kl-partition-refinement}
设 $\mathcal C'$ 细化 $\mathcal C$。记 $\mathcal C'|_C$ 为包含在 $C$ 内的细单元集合，则
\[
 K_{\mathcal C'}-K_{\mathcal C}
 =\sum_{C\in\mathcal C}P_C
   \KL\left(
       \left(\frac{P_{C'}}{P_C}\right)_{C'\in\mathcal C'|_C}
       \middle\Vert
       \left(\frac{Q_{C'}}{Q_C}\right)_{C'\in\mathcal C'|_C}
      \right)\geq0.
\]
正向版本交换两组质量，并将外层权重改为 $Q_C$。因此细化使粗粒度 KL 的标量偏差绝对值和平方偏差单调下降；全部单元都是 singleton 时偏差为零。
\end{proposition}

\begin{proof}
对每个 $C$，利用 $\sum_{C'\subseteq C}P_{C'}=P_C$，将新增贡献写成
\begin{align*}
 \sum_{C'\subseteq C}P_{C'}\log\frac{P_{C'}}{Q_{C'}}
    -P_C\log\frac{P_C}{Q_C}
 &=P_C\sum_{C'\subseteq C}\frac{P_{C'}}{P_C}
       \log\frac{P_{C'}/P_C}{Q_{C'}/Q_C}.
\end{align*}
右边是 $P_C$ 乘以细单元标签之间的条件 KL。对 $C$ 求和即得所述等式。结合粗粒度下界，误差 $K-K_{\mathcal C}$ 非负且随细化下降，其平方也下降。
\end{proof}

\begin{theorem}[粗粒度 KL 的真实梯度及其残差]
\label{thm:kl-partition-gradient}
记 $g_{\mathcal C}=\nabla K_{\mathcal C}$。则
\begin{align}
 g_{\mathcal C}
 &=\sum_CP_C\ell_C\bar s_C
  =\sum_vp(v)\ell_{C(v)}s_v,
 \label{eq:kl-partition-gradient}\\
 g-g_{\mathcal C}
 &=\sum_CP_C\E_{p_C}[(r_A-\ell_C)s_A]
 \label{eq:kl-partition-gradient-residual}\\
 &=\sum_CP_C\left\{
       \operatorname{Cov}_{p_C}(r_A-\ell_C,s_A)
       +D_C\bar s_C\right\}.
 \label{eq:kl-partition-gradient-covariance}
\end{align}
标量与向量的协方差定义为
\[
 \operatorname{Cov}_{p_C}(u_A,s_A)
 =\E_{p_C}[(u_A-\E_{p_C}u_A)(s_A-\bar s_C)].
\]
\end{theorem}

\begin{proof}
由于 $q$ 固定，$\nabla Q_C=0$，而 $\nabla P_C=P_C\bar s_C$。逐项求导得到
\begin{align*}
 \nabla K_{\mathcal C}
 &=\sum_C\{\ell_C\nabla P_C+P_C\nabla\ell_C\}\\
 &=\sum_C\ell_C\nabla P_C+\sum_C\nabla P_C\\
 &=\sum_CP_C\ell_C\bar s_C.
\end{align*}
最后一步用了 $\sum_C\nabla P_C=\nabla1=0$。将 $P_C\bar s_C$ 展开为桶内加权 score，得到第二种表示。完整梯度为 $g=\sum_vp(v)r_vs_v$，与粗粒度梯度逐项相减，得到式~\eqref{eq:kl-partition-gradient-residual}。又因为
\[
 \E_{p_C}[r_A-\ell_C]=D_C,
\]
用 $\E[us]=\operatorname{Cov}(u,s)+\E[u]\E[s]$ 展开每个桶内期望，就得到式~\eqref{eq:kl-partition-gradient-covariance}。
\end{proof}

式~\eqref{eq:kl-partition-gradient-covariance} 把梯度误差分为桶内相关项和条件形状 KL 乘以平均 score 的项。标量粗化误差只包含 $D_C$ 的加权和，梯度误差还依赖 score 的方向与桶内相关性。因此标量下界的单调改进不能直接推出梯度偏差范数单调下降；梯度结论需要单独分析这些向量项。

对正向 KL，冻结权重 $Q_C$ 使粗粒度梯度变成
\[
 g^{\mathrm F}_{\mathcal C}=-\sum_CQ_C\bar s_C,
 \qquad
 g^{\mathrm F}-g^{\mathrm F}_{\mathcal C}
 =-\sum_CQ_C\{\E_{q_C}[s_A]-\bar s_C\}.
\]
这里平均 student score $\bar s_C$ 的积分分布仍是 $p_C$；它不能替换成 teacher 条件均值。第一式来自 $\nabla\ell_C=\bar s_C$，第二式来自完整正向梯度 $-\sum_vq(v)s_v$。

\subsection{一个标量代理与两个同积分的梯度代理}

使用前文的精确积分加残差抽样构造，令标量代理为
\[
 a^{\mathrm K}_v=\ell_{C(v)},\qquad
 I_p(a^{\mathrm K})=K_{\mathcal C}.
\]
在 $A\sim p$ 下，相应无偏数值估计量为
\[
 \widehat K_{\mathcal C}
 =K_{\mathcal C}+r_A-\ell_{C(A)}.
\]
它的均值由链式分解得到完整 $K$，方差为 $\Var_p(r_A-\ell_{C(A)})$。桶内条件形状相同时，残差逐点为零，估计量成为确定性的完整 KL。

梯度积分的目标函数为 $r_vs_v$。同一个粗粒度梯度允许两种可精确积分的代理：
\begin{align*}
 a^{\mathrm{G,const}}_v&=\ell_{C(v)}\bar s_{C(v)},\\
 a^{\mathrm{G,score}}_v&=\ell_{C(v)}s_v.
\end{align*}
两者逐点不同，但精确积分相同：
\[
 I_p(a^{\mathrm{G,const}})
 =\sum_CP_C\ell_C\bar s_C
 =I_p(a^{\mathrm{G,score}})=g_{\mathcal C}.
\]
因此它们产生两个不同的无偏梯度估计量：
\begin{align}
 \widehat g_{\mathcal C}^{\mathrm{const}}
 &=g_{\mathcal C}+r_As_A-\ell_{C(A)}\bar s_{C(A)},
 \label{eq:kl-partition-constant-estimator}\\
 \widehat g_{\mathcal C}^{\mathrm{score}}
 &=g_{\mathcal C}+(r_A-\ell_{C(A)})s_A.
 \label{eq:kl-partition-score-estimator}
\end{align}
两种代理之差为
\[
 \psi_v=\ell_{C(v)}(s_v-\bar s_{C(v)}),
 \qquad \E_{p_C}[\psi_A]=0.
\]
所以改变这两个代理之间的线性组合，只改变桶内零均值波动。它不改变粗粒度精确积分，也不改变每个桶内残差的条件均值。在前文控制矩阵与 Gram 二次型框架中，$\psi_A$ 就是一列已知均值为零的控制特征；其最优系数由残差与该特征的协方差决定，取系数一不自动保证降低方差。

\subsection{头部 singleton 与尾部单桶的特例}

令 $S$ 为冻结的 teacher Top-$J$ 支持，定义 $P=\sum_{v\in S}p(v)$、$Q=\sum_{v\in S}q(v)$、$\tau_p=1-P$、$\tau_q=1-Q$ 和 $J_A=\mathbf1\{A\notin S\}$，本节头尾公式假设 $0<P,Q<1$。头部积分为
\[
 H=\sum_{v\in S}p(v)r_v,\qquad h_g=\sum_{v\in S}p(v)r_vs_v.
\]
令 $\mathcal C=\{\{v\}:v\in S\}\cup\{\mathcal V\setminus S\}$。头部单元满足 $\ell_{\{v\}}=r_v$、$\bar s_{\{v\}}=s_v$；尾部单元满足
\[
 \ell_{\tail}=\log\frac{\tau_p}{\tau_q},\qquad
 \bar s_\tail=\nabla\log\tau_p,\qquad
 b_{\mathrm{TA}}=\ell_{\tail}\bar s_\tail.
\]
因此 $K_{\mathcal C}=H+\tau_p\ell_{\tail}$ 是 Tail-aware KL，且 $g_{\mathcal C}=h_g+\tau_pb_{\mathrm{TA}}$。三个残差构造分别成为
\begin{align*}
 \widehat K_{\mathcal C}
 &=H+\tau_p\ell_{\tail}+J_A(r_A-\ell_{\tail}),\\
 \widehat g_{\mathcal C}^{\mathrm{const}}
 &=h_g+\tau_pb_{\mathrm{TA}}+J_A(r_As_A-b_{\mathrm{TA}}),\\
 \widehat g_{\mathcal C}^{\mathrm{score}}
 &=h_g+\tau_pb_{\mathrm{TA}}+J_A(r_A-\ell_{\tail})s_A.
\end{align*}
第一式是标量残差估计量；第二式是 ResOPD；第三式是 SC-TA。HETS 使用头部精确代理和零尾部代理，因而是另一个代理选择。上述分类从积分目标和代理函数同时确定估计量，不依赖算法名称。

\begin{center}
\begin{tabular}{lll}
\toprule
目标与方法 & 头部近似 $a_v$ & 尾部近似 $a_v$\\
\midrule
数值 HETS & $r_v$ & $0$\\
数值桶残差校正 & $r_v$ & $\ell_{\tail}$\\
梯度 HETS & $r_vs_v$ & $0$\\
梯度 ResOPD & $r_vs_v$ & $b_{\mathrm{TA}}$\\
梯度 SC-TA & $r_vs_v$ & $\ell_{\tail}s_v$\\
\bottomrule
\end{tabular}
\end{center}
表中每一行都使用式~\eqref{eq:kl-master-estimator}；确定积分项是该行的 $I_p(a)$，采样项始终是 $f_A-a_A$。数值 Tail-aware 与梯度粗化则只返回各自行的确定积分项，没有补回残差。

\begin{proposition}[标量精确时，常向量代理仍可有方差]
若 $p_\tail=q_\tail$，则上述标量残差估计量和 SC-TA 都是确定性的完整目标，而 ResOPD 的协方差为
\[
 \operatorname{Cov}(\widehat g^{\mathrm{ResOPD}})
 =\tau_p\ell_{\tail}^2\operatorname{Cov}_{p_\tail}(s_A).
\]
\end{proposition}

\begin{proof}
桶内条件分布相同意味着尾部 $r_v=\ell_{\tail}$。标量残差和 SC-TA 的随机残差于是逐点为零。ResOPD 则仍有随机残差
\[
 J_A\ell_{\tail}(s_A-\bar s_\tail).
\]
它在尾部条件下均值为零，在头部为零，因此总体均值为零，二阶矩恰好是命题中的协方差。
\end{proof}

例如取 $p=(0.4,0.2,0.4)$、$q=(0.7,0.1,0.2)$ 和 $S=\{1\}$，则两个尾部条件分布均为 $(1/3,2/3)$，且 $\tau_p=0.6$、$\ell_{\tail}=\log2$。在完整 softmax logits 参数化下，$s_v=e_v-p$，尾部 score 的协方差迹为 $1-(1/3)^2-(2/3)^2=4/9$。ResOPD 的方差因此为 $\frac4{15}(\log2)^2>0$；标量残差估计量与 SC-TA 的方差均为零。这个例子直接区分了可精确积分的 KL 数值、粗粒度真实梯度以及不同梯度代理留下的抽样波动。

\section{KL 数值：从确定性代理到无偏残差校正}
\label{sec:kl-scalar}
本节固定目标 $K=I_p(r)$，所以总定理中的 $f_v=r_v$ 是标量。随机估计量默认使用 $\rho=p$；任意 $\rho$ 的权重与支持条件见第~\ref{sec:kl-sampling} 节。先给出 Top-$J$ 记号，再分类确定性代理、全词表单样本控制和稀疏头尾控制。每个估计量的误差都相对同一个完整 $K$。
\subsection{头部与尾部记号}

沿用前节的 $S,P,Q,\tau_p,\tau_q,J_A,H$ 和 $\ell_{\tail}$。令
\[
 K=\KL(p\Vert q)=\sum_vp(v)r_v,
 \qquad H=\sum_{v\in S}p(v)r_v,
 \]
并令尾部条件分布为
\[
 p_\tail(v)=\frac{p(v)}{\tau_p},\qquad
 q_\tail(v)=\frac{q(v)}{\tau_q},\qquad v\notin S.
\]
定义尾部质量项和尾部形状项
\[
 M=\tau_p\log\frac{\tau_p}{\tau_q},\qquad
 B=\tau_p\KL(p_\tail\Vert q_\tail)\geq0.
\]

\begin{corollary}[头尾分区的 KL 链式分解]
\label{thm:chapter3-kl-chain}
在 $p\ll q$ 且 $0<\tau_p,\tau_q<1$ 时，
\begin{equation}
 K=H+M+B.
\label{eq:chapter3-kl-chain}
\end{equation}
\end{corollary}

\begin{proof}
在定理~\ref{thm:kl-partition-chain} 中取头部 singleton 与单个尾桶。singleton 的条件 KL 为零，粗粒度 KL 是 $H+M$，唯一剩余的条件形状项是 $B$，故得到所述恒等式。
\end{proof}

该分解说明 sparse teacher payload 能够知道 $H$、$P$、$Q$，因而知道 $M$，但不知道尾部形状项 $B$。所有 Top-$J$ scalar estimator 的差异都可以直接归结为是否丢弃、聚合或抽样估计 $B$。

\subsection{确定性估计量}

截断 Top-$J$ 估计量定义为
\[
 \widehat K^{\mathrm{trunc}}=H.
\]
由式~\eqref{eq:chapter3-kl-chain}，其条件偏差和方差为
\begin{equation}
 \operatorname{Bias}(\widehat K^{\mathrm{trunc}})=-M-B,
 \qquad
 \Var(\widehat K^{\mathrm{trunc}})=0.
 \label{eq:chapter3-trunc-bias-var}
\end{equation}
它删除了尾部贡献；偏差 $-M-B$ 可正可负，因为 $M$ 的符号可以正也可以负，而 $B$ 总是非负。

Tail-aware scalar 估计量把尾部聚合为一个伪 token：
\[
 \widehat K^{\mathrm{TA}}=H+M.
\]

\begin{proposition}[Tail-aware scalar 的偏差]
\label{prop:chapter3-ta-scalar}
\[
 \operatorname{Bias}(\widehat K^{\mathrm{TA}})=-B,
 \qquad
 \Var(\widehat K^{\mathrm{TA}})=0.
\]
因此它是完整 KL 的确定性下界，且当且仅当尾部条件分布相等时达到完整 KL。
\end{proposition}

\begin{proof}
将 $\widehat K^{\mathrm{TA}}=H+M$ 与 $K=H+M+B$ 相减即可。$B\geq0$ 来自 KL 非负性。
\end{proof}

Normalized Top-$J$ 不是完整 KL 的估计量，而是头部条件分布的另一个目标：
\[
 K^{\mathrm{norm}}_J
 =\KL\left(\frac{p|_S}{P}\middle\|\frac{q|_S}{Q}\right)
 =\frac{H}{P}+\log\frac{Q}{P}.
\]
它的条件方差为零，但相对完整 KL 的确定性偏差为
\begin{equation}
 \operatorname{Bias}(K^{\mathrm{norm}}_J)
 =\frac{H}{P}+\log\frac{Q}{P}-(H+M+B).
 \label{eq:chapter3-normalized-bias}
\end{equation}
这个偏差没有固定符号，因为 Normalized Top-$J$ 同时改变了头部权重并删除了尾部形状。

\subsection{单 token KL 估计器的基准族}

先记录不依赖 Top-$J$ 的单 token 族。令 $A\sim p$，$r=r_A$，并假设相应矩存在。反向 KL 的直接 estimator 为
\[
 \widehat K_1=r.
\]
它满足
\[
 \E[\widehat K_1]=K,\qquad
 \operatorname{Bias}(\widehat K_1)=0,\qquad
 \Var(\widehat K_1)=\Var_p(r).
\]
常用的平方 estimator 为
\[
 \widehat K_2=\frac12r^2,
\]
其相对于 $K$ 的偏差和方差为
\begin{equation}
 \operatorname{Bias}(\widehat K_2)=\frac12\E_p[r^2]-K,
 \qquad
 \Var(\widehat K_2)=\frac14\left(\E_p[r^4]-\E_p[r^2]^2\right).
 \label{eq:chapter3-k2-bias-var}
\end{equation}
非负修正 estimator 为
\[
 \widehat K_3=e^{-r}+r-1.
\]
因为 $e^{-r}=q(A)/p(A)$，有
\[
 \E_p[e^{-r}]=\sum_vq(v)=1,
\]
故 $\E[\widehat K_3]=K$，且
\begin{equation}
 \Var(\widehat K_3)=\Var_p(e^{-r}+r-1).
 \label{eq:chapter3-k3-var}
\end{equation}
K3 的方差还可展开为
\[
 \Var(\widehat K_3)
 =\Var_p(r)+\Var_p(e^{-r})+2\operatorname{Cov}_p(r,e^{-r}).
\]
因此 K3 比 K1 方差小的充要条件是
\[
 \Var_p(e^{-r})+2\operatorname{Cov}_p(r,e^{-r})<0.
\]
K3 的逐样本非负性来自 $e^{-r}\geq1-r$，并不推出全局方差排序。对 $q(v)=p(v)(1+\varepsilon u_v)$、$\E_p[u]=0$ 的局部扰动，逐项 Taylor 展开给出
\[
 r_v=-\varepsilon u_v+\frac{\varepsilon^2u_v^2}{2}+O(\varepsilon^3),
 \quad K=\frac{\varepsilon^2}{2}\E_p[u^2]+O(\varepsilon^3),
\]
\[
 \widehat K_2=\frac{\varepsilon^2u_A^2}{2}+O(\varepsilon^3),
 \quad \widehat K_3=\frac{\varepsilon^2u_A^2}{2}+O(\varepsilon^3).
\]
这解释了 K2 的局部二阶正确性，以及 K3 在两分布接近时将单点波动从一阶降到二阶的机制。
K1 是总定理中 $a=0$ 的特例。K3 则取 $a_v=1-e^{-r_v}$，因为 $I_p(a)=1-\sum_vq(v)=0$，总定理给出 $r_A-a_A=e^{-r_A}+r_A-1$。K2 改变了数值目标而未补回残差，所以需要直接报告上面的偏差。

\subsection{采样尾部的无偏 scalar 估计量}

令
\[
 \mu_r=\E[r_A\mid A\notin S]
 =\log\frac{\tau_p}{\tau_q}+\KL(p_\tail\Vert q_\tail).
\]
则 $\tau_p\mu_r=M+B$。头部精确、尾部采样（Head-Exact Tail-Sampled，HETS）的数值估计量为
\begin{equation}
 \widehat K^{\mathrm{HETS}}=H+J_A r_A.
 \label{eq:chapter3-ema-scalar}
\end{equation}
它是总定理中头部 $a_v=r_v$、尾部 $a_v=0$ 的特例。本章后续统一称为 HETS 数值估计量；EMA 只在对应原项目名称时出现。时间上的指数滑动平均不属于此命名。

\begin{theorem}[HETS 数值的无偏性与方差]
\label{thm:chapter3-ema-scalar}
若 $A\sim p$ 且 $\E_p[r_A^2]<\infty$，则
\[
 \E[\widehat K^{\mathrm{HETS}}]=K,
\]
并且
\begin{align}
 \Var(\widehat K^{\mathrm{HETS}})
 &=\tau_p\Var(r_A\mid A\notin S)+P\tau_p\mu_r^2\\
 &=\sum_{v\notin S}p(v)r_v^2
   -\left(\sum_{v\notin S}p(v)r_v\right)^2.
 \label{eq:chapter3-ema-scalar-var}
\end{align}
\end{theorem}

\begin{proof}
尾部修正的期望为
\[
 \E[J_A r_A]=\sum_{v\notin S}p(v)r_v=\tau_p\mu_r=M+B.
\]
加上 $H$ 得到 $K$。对随机变量 $J_A r_A$ 按 $J_A$ 使用全方差公式：头部条件下它为零，尾部条件下它为 $r_A$，因此
\[
 \Var(J_A r_A)=\tau_p\Var(r_A\mid A\notin S)+P\tau_p\mu_r^2.
\]
\end{proof}

还可以用尾部桶 scalar
\[
 \ell_{\tail}=\log\frac{\tau_p}{\tau_q}
\]
作残差化：
\begin{equation}
 \widehat K^{\mathrm{Res\text{-}scalar}}
 =H+\tau_p\ell_{\tail}+J_A(r_A-\ell_{\tail}).
 \label{eq:chapter3-res-scalar}
\end{equation}
它满足
\[
 \E[\widehat K^{\mathrm{Res\text{-}scalar}}]=K,
\]
且
\begin{equation}
 \Var(\widehat K^{\mathrm{Res\text{-}scalar}})
 =\tau_p\Var(r_A\mid A\notin S)+P\tau_p(\mu_r-\ell_{\tail})^2.
 \label{eq:chapter3-res-scalar-var}
\end{equation}
所以 scalar 残差化相对 HETS 的方差差为
\begin{equation}
 \Var(\widehat K^{\mathrm{HETS}})-\Var(\widehat K^{\mathrm{Res\text{-}scalar}})
 =P\tau_p(2\mu_r\ell_{\tail}-\ell_{\tail}^2).
 \label{eq:chapter3-res-scalar-gap}
\end{equation}
它严格降低方差当且仅当 $|\mu_r-\ell_{\tail}|<|\mu_r|$。当固定的 $\ell_{\tail}\ne0$ 且尾部内部形状 KL 足够小时，$\mu_r$ 足够接近 $\ell_{\tail}$，严格降方差条件成立；若 $\ell_{\tail}=0$，两估计量逐样本相同。

\begin{proposition}[标量尾部控制族]
对任意在抽样前固定的实数 $\beta$，令
\[
 \widehat K(\beta)=H+\tau_p\beta+J_A(r_A-\beta).
\]
则它对 $K$ 无偏，且
\[
 \Var(\widehat K(\beta))
 =\tau_p\Var_{\tail}(r_A)+P\tau_p(\mu_r-\beta)^2.
\]
因此最优标量控制为 $\beta^*=\mu_r$；HETS 取 $\beta=0$，式~\eqref{eq:chapter3-res-scalar} 取 $\beta=\ell_{\tail}$。
\end{proposition}

\begin{proof}
先计算均值：
\[
 \E[\widehat K(\beta)]=H+\tau_p\beta+\tau_p(\mu_r-\beta)=H+\tau_p\mu_r=K.
\]
再计算中心化随机项的二阶矩：
\begin{align*}
 \Var(\widehat K(\beta))
 &=\E[J_A(r_A-\beta)^2]-\{\E[J_A(r_A-\beta)]\}^2\\
 &=\tau_p\{\Var_{\tail}(r_A)+(\mu_r-\beta)^2\}
   -\tau_p^2(\mu_r-\beta)^2\\
 &=\tau_p\Var_{\tail}(r_A)+P\tau_p(\mu_r-\beta)^2.
\end{align*}
最后一项非负并在 $\beta=\mu_r$ 时为零，给出最优值。将 $\beta=0$ 与 $\beta=\ell_{\tail}$ 代入并相减即得式~\eqref{eq:chapter3-res-scalar-gap}。
\end{proof}

这个 scalar 构造是本章从控制族导出的 KL 数值估计量，应与论文中以向量梯度为目标的 ResOPD 算法区别命名。因为 $\mu_r=\ell_{\tail}+D_{\tail}$，其中 $D_{\tail}=\KL(p_\tail\Vert q_\tail)$，方差改变量也等于
\[
 P\tau_p\{\ell_{\tail}^2+2\ell_{\tail}D_{\tail}\}.
\]
当 $\ell_{\tail}\geq0$ 时它非负；当 $\ell_{\tail}<0$ 时，它在 $D_{\tail}>-\ell_{\tail}/2$ 下为负。这给出 scalar 残差化反而增大方差的明确条件。

\section{KL 梯度：相同积分下的不同控制}
\label{sec:kl-vector}
本节固定目标 $g=I_p(rs)$，总定理中的 $f_v=r_vs_v$ 是向量，随机估计量默认使用 $\rho=p$。数值估计量的分类逐项迁移到这个目标，但其控制函数和最优系数须在向量二阶矩下重新计算。沿用前节的 $h_g,\bar s_\tail,b_{\mathrm{TA}}$，并定义
\[
 \mu_g=\E_p[r_As_A\mid J_A=1],\qquad
 d_P=\nabla P=\sum_{v\in S}p(v)s_v,\qquad g=h_g+\tau_p\mu_g.
\]
头部截断梯度为 $h_g$，其偏差和方差为
\begin{equation}
 \operatorname{Bias}(h_g)=-\tau_p\mu_g,
 \qquad \Var(h_g)=0.
 \label{eq:chapter3-head-grad-bias}
\end{equation}

\subsection{HETS 与 ResOPD 的向量残差化}

尾部残差化在线策略蒸馏（On-Policy Distillation with Tail Residualization，ResOPD）从 HETS 的尾部梯度中减去可计算控制向量，再解析加回同一控制的期望。

HETS 梯度 estimator 定义为
\begin{equation}
 \widehat g^{\mathrm{HETS}}=h_g+J_A r_As_A.
 \label{eq:chapter3-hets-grad}
\end{equation}
它无偏，且
\begin{equation}
 \Var(\widehat g^{\mathrm{HETS}})
 =\tau_p\Var_{\tail}(r_As_A)+P\tau_p\lVert\mu_g\rVert_2^2.
 \label{eq:chapter3-hets-var}
\end{equation}

将尾部聚合为伪 token，定义
\begin{equation}
 b_{\mathrm{TA}}=\log\frac{\tau_p}{\tau_q}\,\nabla_\theta\log\tau_p.
 \label{eq:chapter3-resopd-control}
\end{equation}
这个控制向量来自定理~\ref{thm:kl-partition-gradient} 的粗粒度梯度。对头尾分区直接代入该定理，得到 $\nabla(H+\tau_p\ell_{\tail})=h_g+\tau_pb_{\mathrm{TA}}$。因此完整梯度与粗粒度梯度之差为 $\tau_p(\mu_g-b_{\mathrm{TA}})$。使用 $J_A(r_As_A-b_{\mathrm{TA}})$ 对该残差作无偏抽样，再加上粗粒度精确梯度，便得到 ResOPD。
对任意 action-independent 向量 $b$（与标量尾部形状项 $B$ 区分），定义统一控制族
\begin{equation}
 \widehat g(b)=h_g+\tau_pb+J_A(r_As_A-b).
 \label{eq:chapter3-control-family}
\end{equation}

\begin{theorem}[ResOPD 控制族的无偏性与方差]
\label{thm:chapter3-resopd}
在 $A\sim p$ 且 $b$ 在抽样前冻结时，
\[
 \E[\widehat g(b)]=g,
\]
并且
\begin{equation}
 \Var(\widehat g(b))
 =\tau_p\Var_{\tail}(r_As_A)
  +P\tau_p\lVert\mu_g-b\rVert_2^2.
 \label{eq:chapter3-control-var}
\end{equation}
HETS 是 $b=0$；ResOPD 是 $b=b_{\mathrm{TA}}$。
\end{theorem}

\begin{proof}
条件期望为
\[
 \E[\widehat g(b)]
 =h_g+\tau_pb+\tau_p(\mu_g-b)=g.
\]
给定 $J_A=0$ 时 estimator 为 $h_g+\tau_pb$，给定 $J_A=1$ 时其条件均值为 $h_g+\tau_pb+\mu_g-b$，两者均值差为 $\mu_g-b$。尾部条件下的随机部分为 $r_As_A-b$，其协方差与 $r_As_A$ 相同。对事件 $J_A$ 使用全协方差公式，得到式~\eqref{eq:chapter3-control-var}。
\end{proof}

ResOPD 与 HETS 的方差差为
\begin{equation}
 \Var(\widehat g^{\mathrm{HETS}})-\Var(\widehat g^{\mathrm{ResOPD}})
 =P\tau_p\left(2\langle\mu_g,b_{\mathrm{TA}}\rangle-\lVert b_{\mathrm{TA}}\rVert_2^2\right).
 \label{eq:chapter3-hets-resopd-gap}
\end{equation}
因此 ResOPD 的优势条件是
$\lVert\mu_g-b_{\mathrm{TA}}\rVert_2<\lVert\mu_g\rVert_2$。这正是 ResOPD 与 scalar KL 残差化的共同结构，区别只在于 scalar 的尾部均值是 $\mu_r$，向量梯度的尾部均值是 $\mu_g$。

该方差结论也有完整的协方差版本。记 $\Sigma_{\tail}=\operatorname{Cov}_{\tail}(r_As_A)$，则
\[
 \operatorname{Cov}(\widehat g(b))
 =\tau_p\Sigma_{\tail}
  +P\tau_p(\mu_g-b)(\mu_g-b)^\top.
\]
取迹即得到前面的方差公式。最优的无约束向量控制是 $b^*=\mu_g$，此时只剩 $\tau_p\Sigma_{\tail}$；该向量需要未知尾部，因此理论最优值不能作为 sparse payload 下的可计算算法。ResOPD 使用可计算的 $b_{\mathrm{TA}}$ 逼近它。

HETS 也不保证对每个分布都比直接 sampled-token 方差小。令 $\mu_{g,H}=\E[r_As_A\mid A\in S]$、$V_H=\Var(r_As_A\mid A\in S)$，全方差公式给出
\begin{align*}
 \Var(r_As_A)&=PV_H+\tau_p\Var_{\tail}(r_As_A)
 +P\tau_p\lVert\mu_{g,H}-\mu_g\rVert^2,\\
 \Var(r_As_A)-\Var(\widehat g^{\mathrm{HETS}})
 &=PV_H+P\tau_p\{\lVert\mu_{g,H}\rVert^2
 -2\langle\mu_{g,H},\mu_g\rangle\}.
\end{align*}
头部精确求和删除了头部内部方差，但同时改变了事件间均值差，因此应使用这个差值判断排序。

\subsection{多样本、均方误差与支持集边界}

以上随机 estimator 均按单个动作计算。对固定 prefix 下 $n$ 个独立动作，使用同一确定性头部和控制项、对随机修正取平均，则偏差保持不变、协方差除以 $n$。标量均方误差为
\[
 \operatorname{MSE}(\widehat K)=\operatorname{Bias}(\widehat K)^2+\Var(\widehat K),
\]
向量均方误差为
\[
 \E\lVert\widehat g-g\rVert^2
 =\lVert\operatorname{Bias}(\widehat g)\rVert^2+\Var(\widehat g).
\]
因此确定性截断的零抽样方差不等于零估计误差；增加样本也无法消除它的偏差。这里的 $1/n$ 结论要求条件独立，同一条 trajectory 上不同 prefix 的 token 不能直接视为同一固定分布的独立重复。

当 $S=\mathcal V$ 时，所有尾部修正按连续延拓取零，头部求和等于完整 KL 或其完整 score 梯度，方差为零；当 $S=\varnothing$ 时，HETS 退化为直接 sampled-token，$\tau_p=\tau_q=1$ 使 ResOPD 控制为零。存在零概率时必须先检查绝对连续性及矩条件；尤其 K3 的 $\E_p[q(A)/p(A)]=1$ 要求 $q$ 的质量全部落在 $p$ 的支持上。

\section{双控制族：事件间方差与尾部内部方差}
\label{sec:kl-two-controls}

单控制 ResOPD 只改变头部与尾部两个事件的条件均值差。本节将总定理中的近似扩展为一个尾部常向量与一个条件零均值方向，并使用同一套 $r_v,s_v,J_A,\tau_p,\mu_g,b_{\mathrm{TA}}$ 记号。定义尾部平均 score
\[
 \bar s_{\tail}=\E_p[s_A\mid J_A=1]
 =\frac{1}{\tau_p}\sum_{v\notin S}p(v)s_v
 =\nabla_\theta\log\tau_p.
\]
沿用 $\ell_{\tail}=\log(\tau_p/\tau_q)$，于是 $b_{\mathrm{TA}}=\ell_{\tail}\bar s_{\tail}$。定义向量
\[
 \psi_v=\ell_{\tail}(s_v-\bar s_{\tail}),\qquad
 R_v^{\mathrm{shape}}=(r_v-\ell_{\tail})s_v,\qquad v\notin S.
\]
其中 $\psi_v$ 是尾部条件零均值波动，$R_v^{\mathrm{shape}}$ 是尾部形状残差。逐点恒等式为
\begin{equation}
 r_vs_v=b_{\mathrm{TA}}+\psi_v+R_v^{\mathrm{shape}},\qquad
 \E_p[\psi_A\mid J_A=1]=0.
 \label{eq:kl-tail-fine-decomposition}
\end{equation}
这些定义没有改变目标或采样分布。

\subsection{双控制是总定理的两个正交方向}

在式~\eqref{eq:kl-master-estimator} 中取头部 $a_v=r_vs_v$、尾部 $a_v=b+\gamma\psi_v$。因为 $\E_{p_\tail}\psi_A=0$，精确积分为 $h_g+\tau_pb$，立即得到下面的双控制估计量。

其正交结构也可以直接从 Gram 矩阵看出。尾部常向量的第 $j$ 个坐标对应零均值控制 $(J_A-\tau_p)e_j$，尾部内部控制为 $J_A\psi_A$，其中 $e_j$ 为参数空间的标准基向量。两列的 Gram 交叉项为
\[
 \E_p\langle(J_A-\tau_p)e_j,J_A\psi_A\rangle
 =P\tau_p\E_{p_\tail}\langle e_j,\psi_A\rangle=0.
\]
因此 on-policy 下的 Gram 矩阵分成事件控制和尾部内部控制两个块。后面的两个平方代价来自这一块对角结构，而非两个控制之间的巧合抵消。

\begin{theorem}[双控制估计量及其精确协方差]
\label{thm:kl-two-controls}
对抽样前固定的向量 $b$ 与标量 $\gamma$，定义
\[
 \widehat g(b,\gamma)
 =h_g+\tau_pb+J_A(r_As_A-b-\gamma \psi_A),\qquad A\sim p.
\]
则 $\E_p[\widehat g(b,\gamma)]=g$，并且
\begin{align}
 \operatorname{Cov}(\widehat g(b,\gamma))
 &=\tau_p\operatorname{Cov}_{\tail}(r_As_A-\gamma \psi_A)
   +P\tau_p(\mu_g-b)(\mu_g-b)^\top, \label{eq:kl-two-control-cov}\\
 \Var(\widehat g(b,\gamma))
 &=\tau_p\Var_{\tail}(r_As_A-\gamma \psi_A)
   +P\tau_p\lVert\mu_g-b\rVert^2. \label{eq:kl-two-control-var}
\end{align}
\end{theorem}

\begin{proof}
尾部条件下 $\psi_A$ 的均值为零，所以
\[
 \E_p[J_A(r_As_A-b-\gamma \psi_A)]
 =\tau_p(\mu_g-b).
\]
加上确定性部分 $h_g+\tau_pb$ 就是 $g$。给定 $J_A=0$ 时，估计量是常向量 $h_g+\tau_pb$；给定 $J_A=1$ 时，均值为 $h_g+\tau_pb+\mu_g-b$，协方差为 $\operatorname{Cov}_{\tail}(r_As_A-\gamma \psi_A)$。两分支均值差为 $\mu_g-b$。将这两个条件协方差与条件均值协方差相加，得到式~\eqref{eq:kl-two-control-cov}；取迹得到式~\eqref{eq:kl-two-control-var}。
\end{proof}

\begin{theorem}[双控制的最优值与精确代价]
令
\[
 V_\psi=\Var_{\tail}(\psi_A),\qquad
 C_\psi=\E_{\tail}\langle r_As_A-\mu_g,\psi_A\rangle.
\]
当 $V_\psi>0$ 时，迹方差最优控制为
\[
 b^*=\mu_g,\qquad \gamma^*=\frac{C_\psi}{V_\psi}.
\]
任意 $b,\gamma$ 相对最优值的方差差恰好为
\begin{equation}
 \Var(\widehat g(b,\gamma))-\Var(\widehat g(b^*,\gamma^*))
 =P\tau_p\lVert b-\mu_g\rVert^2
  +\tau_pV_\psi(\gamma-\gamma^*)^2.
 \label{eq:kl-two-control-regret}
\end{equation}
当 $V_\psi=0$ 时，$\gamma$ 不影响估计量的方差。
\end{theorem}

\begin{proof}
由 $\E_{\tail}\psi_A=0$，
\[
 \Var_{\tail}(r_As_A-\gamma \psi_A)
 =\Var_{\tail}(r_As_A)-2\gamma C_\psi+\gamma^2V_\psi.
\]
对 $\gamma$ 配方，得到
\[
 \Var_{\tail}(r_As_A-\gamma \psi_A)
 =\Var_{\tail}(r_As_A)-\frac{C_\psi^2}{V_\psi}
  +V_\psi(\gamma-\gamma^*)^2.
\]
再代入式~\eqref{eq:kl-two-control-var}，两个非负平方项分别由 $b=\mu_g$ 与 $\gamma=\gamma^*$ 消去。若 $V_\psi=0$，则 $\psi_A=0$ 在尾部条件分布下几乎处处成立。
\end{proof}

\subsection{HETS、ResOPD 与 SC-TA 的统一分类}

样本修正 Tail-aware（Sample-Corrected Tail-Aware，SC-TA）与 ResOPD 的区别在尾部内部控制。双控制族中的坐标为
\begin{center}
\begin{tabular}{lcc}
\toprule
估计量 & $b$ & $\gamma$\\
\midrule
HETS & $0$ & $0$\\
ResOPD & $b_{\mathrm{TA}}$ & $0$\\
SC-TA & $b_{\mathrm{TA}}$ & $1$\\
不可计算的联合最优值 & $\mu_g$ & $\gamma^*$\\
\bottomrule
\end{tabular}
\end{center}
三种可计算估计量在 $A\sim p$ 下均无偏。令 $g_{\mathrm{TA}}=h_g+\tau_pb_{\mathrm{TA}}$，则逐样本有
\begin{align*}
 \widehat g^{\mathrm{ResOPD}}&=g_{\mathrm{TA}}+J_A(R_A^{\mathrm{shape}}+\psi_A),\\
 \widehat g^{\mathrm{SC\text{-}TA}}&=g_{\mathrm{TA}}+J_AR_A^{\mathrm{shape}},\\
 \widehat g^{\mathrm{SC\text{-}TA}}-\widehat g^{\mathrm{ResOPD}}&=-J_A\psi_A.
\end{align*}
这个差的期望为零，逐样本一般不为零。

\begin{proposition}[SC-TA 与 ResOPD 的方差排序]
当 $V_\psi>0$ 时，
\[
 \Var(\widehat g^{\mathrm{ResOPD}})
 -\Var(\widehat g^{\mathrm{SC\text{-}TA}})
 =\tau_p(2C_\psi-V_\psi)
 =\tau_pV_\psi(2\gamma^*-1).
\]
因此 SC-TA 方差更小当且仅当 $\gamma^*>1/2$。一个充分条件为
\[
 \Var_{\tail}(R_A^{\mathrm{shape}})<\frac14\Var_{\tail}(\psi_A).
\]
\end{proposition}

\begin{proof}
在同一个 $b=b_{\mathrm{TA}}$ 下，事件间方差相同。将 $\gamma=0,1$ 代入尾部内部二次式并相减，得到第一式。由 $r_As_A=b_{\mathrm{TA}}+R_A^{\mathrm{shape}}+\psi_A$，
\[
 C_\psi=V_\psi+\E_{\tail}\langle R_A^{\mathrm{shape}}-\E_{\tail}R_A^{\mathrm{shape}},\psi_A\rangle.
\]
Cauchy--Schwarz 不等式给出交叉项的下界
\[
 \E_{\tail}\langle R_A^{\mathrm{shape}}-\E_{\tail}R_A^{\mathrm{shape}},\psi_A\rangle
 \geq-\sqrt{\Var_{\tail}(R_A^{\mathrm{shape}})V_\psi}.
\]
给定充分条件时，右侧严格大于 $-V_\psi/2$，故 $2C_\psi-V_\psi>0$。
\end{proof}

当尾部条件形状完全相同，$r_v=\ell_{\tail}$ 对所有尾部 token 成立，此时 $R_A^{\mathrm{shape}}=0$，SC-TA 等于确定性的完整梯度；ResOPD 仍保留 $J_A\psi_A$ 的零均值波动。因此即使尾部桶已经完全描述 KL 数值，ResOPD 的梯度估计也可以有非零方差。若 $V_\psi=0$，两者逐样本相同。这是区分标量 KL 与向量梯度的一个直接例子。

上述所有排序针对协方差迹。较小的迹不推出协方差矩阵在半正定序下更小；事件间协方差之差是两个秩一矩阵之差，某些方向上的方差仍可增大。

当动作改由 $\rho$ 抽样时，令 $w_A=p(A)/\rho(A)$，并假设 $\rho$ 覆盖全部尾部。正确的重要性控制变成 $(J_Aw_A-\tau_p)e_j$ 与 $J_Aw_A\psi_A$。其交叉项为
\[
 \sum_{v\notin S}\frac{p(v)^2}{\rho(v)}\langle e_j,\psi_v\rangle,
\]
一般不为零：$\psi$ 的零均值只针对 $p_\tail$，而这个交叉矩使用 $p^2/\rho$。因此非 on-policy 的最优常向量与内部控制应由完整 Gram 矩阵联合求解。仅在块交叉项确实为零时，才可以分别优化两个控制。

\section{自动微分：数值与梯度的分别实现}
\label{sec:kl-autodiff}
前面的统一框架估计已经指定的积分。本节检查标量程序的反传究竟对应哪个积分。若一个程序同时用于 KL 日志与训练，审计对象应记录为 forward 数值和 backward 向量组成的有序对。
本节涉及随机梯度的期望与偏差时默认 $\rho=p$；固定索引的逐样本求导规则本身不依赖抽样分布。

\subsection{一个函数的两种求导路径}

\begin{theorem}[固定样本反传与采样测度补偿]
\label{thm:kl-autodiff-transform}
令 $\phi$ 是不另依赖参数 $\theta$ 的可微实函数，$A\sim p$，反传时固定样本索引。定义
\[
 F_A=\phi(r_A),\qquad
 \kappa_A=\frac{p(A)}{\sg(p(A))},\qquad
 L_A=\kappa_A\phi(r_A).
\]
两者的 forward 值逐样本相同，固定样本导数分别为
\[
 \nabla F_A=\phi'(r_A)s_A,\qquad
 \nabla L_A=\{\phi(r_A)+\phi'(r_A)\}s_A.
\]
后者的期望等于 $\nabla I_p(\phi(r))$。若仅将系数停止梯度，$\kappa_A\sg(\phi(r_A))$ 的导数则为 $\phi(r_A)s_A$。
\end{theorem}

\begin{proof}
固定索引下 $\nabla r_A=s_A$，链式法则给出第一式。forward 时 $\kappa_A=1$；分母停止梯度，因而 $\nabla\kappa_A=s_A$。乘积法则给出第二式。再由式~\eqref{eq:kl-measure-derivative}，
\[
 \nabla I_p(\phi(r))
 =I_p\bigl(\phi'(r)s+\phi(r)s\bigr),
\]
等于第二种导数的期望。停止整个 $\phi$ 时只保留 $\nabla\kappa_A$ 一项，得到最后的规则。
\end{proof}

这个变换把 K1--K4 的差异归到同一个问题：给定 forward 函数后，程序保留了哪些乘积法则项。数值无偏是 $I_p(\phi(r))=K$，裸反传无偏则是 $I_p(\phi'(r)s)=g$，两者是不同等式。

\begin{center}
\begin{tabular}{lll}
\toprule
程序 & forward 值 & 固定样本 backward\\
\midrule
K1 & $r_A$ & $s_A$\\
K2 & $r_A^2/2$ & $r_As_A$\\
K3 & $e^{-r_A}+r_A-1$ & $(1-e^{-r_A})s_A$\\
K4 & $r_A$ & $r_As_A$\\
K4$^{\mathrm{live}}$ & $r_A$ & $(r_A+1)s_A$\\
K3$^{++}$ & $e^{-r_A}+r_A-1$ & $r_As_A$\\
\bottomrule
\end{tabular}
\end{center}
数值方差按第二列计算，梯度方差按第三列计算。K5 的目标是正向 KL，随后单独代入同一求导规则。

\subsection{K4 与固定样本的梯度替换}
K1 与 K3 的数值都无偏；K2 是以二阶矩换取局部稳定性的有偏 proxy。若需要同时让数值和自动微分梯度对应同一个 $K$，使用停止 log-ratio 系数的 K4：
\[
 \kappa_A=\frac{p(A)}{\operatorname{sg}(p(A))},\qquad
 \widehat K_4=\kappa_A\sg(r_A).
\]
数值上 $\kappa_A=1$，所以 $\E[\widehat K_4]=K$；反传时
\[
 \nabla_\theta\widehat K_4=r_As_A,
\]
因此
\[
 \E_p[\nabla_\theta\widehat K_4]=\nabla_\theta K=g.
\]
K4 与 K1 的 forward 值逐样本相同，故 K4 的数值偏差为零，数值方差也为 $\Var_p(r)$。stop-gradient 只改变求导规则。另定义本项目单 token 测试采用的 live-log-ratio 变体
\[
 \widehat K_4^{\mathrm{live}}=\kappa_A r_A,\qquad
 \nabla\widehat K_4^{\mathrm{live}}=(r_A+1)s_A.
\]
它的数值仍与 K4 相同，梯度期望也相同，但梯度方差不同。后文始终保留上标 $\mathrm{live}$，不将两种反传路径共用一个名称。

\begin{theorem}[K4 的数值与梯度无偏性]
\label{thm:k4}
在本章正支持与固定参考假设下，$\widehat K_4$ 和 $\widehat K_4^{\mathrm{live}}$ 的数值均对 $K$ 无偏，自动微分梯度均对 $g$ 无偏。
\end{theorem}
\begin{proof}
两者的数值为 $r_A$，期望是 $\sum_vp(v)r_v=K$。两者的梯度分别为 $r_As_A$ 与 $(r_A+1)s_A$，期望之差为 $\sum_vp(v)s_v=0$，因此都等于 $g$。
\end{proof}

\subsection{K3$^{++}$ 与正向 KL 的 K5}

EMA-PG 的公开实现使用上述停止系数 K4，并定义 K3$^{++}$ 与 K5；以下公式先去掉 clipping，单独分析数学对象。\footnote{代码版本固定为 EMA-PG \texttt{f9b9776}，定义见该版本的 \href{https://github.com/LunjunZhang/ema-pg/blob/f9b977632ad948ce935432a517c22291bfaba562/search/verl/trainer/ppo/core_algos.py}{\texttt{core\_algos.py}}。}
\[
 \widehat K_{3^{++}}=\kappa_A(e^{-r_A}+r_A-1).
\]
它的数值与 K3 相同，因此对反向 $K$ 无偏，数值方差也与 K3 相同。其反传为
\[
 \nabla\widehat K_{3^{++}}
 =\{e^{-r_A}+r_A-1+1-e^{-r_A}\}s_A=r_As_A.
\]
因此它与 K4 的 backward 逐样本相同，梯度偏差为零，梯度方差也完全相同；数值方差的差异不能解释成梯度方差的差异。

K5 对应正向目标 $K^{\mathrm F}=\KL(q\Vert p)$，定义为
\[
 \widehat K_5=-\sg(e^{-r_A})r_A+\log p(A)-\sg(\log p(A)).
\]
其 forward 值为 $-e^{-r_A}r_A$。在 $A\sim p$ 下，
\[
 \E_p[\widehat K_5]=-\sum_vq(v)r_v=K^{\mathrm F},\qquad
 \Var(\widehat K_5)=\E_p[e^{-2r_A}r_A^2]-(K^{\mathrm F})^2.
\]
它相对本章默认反向目标 $K$ 的数值偏差是 $K^{\mathrm F}-K$，不能与 K1--K4 按同一个无偏目标直接并列。反传为
\[
 \nabla\widehat K_5=(1-e^{-r_A})s_A,\qquad
 \E_p[\nabla\widehat K_5]=-\sum_vq(v)s_v=\nabla K^{\mathrm F}.
\]
其梯度方差为
\[
 \E_p[(1-e^{-r_A})^2\lVert s_A\rVert^2]
 -\left\lVert\sum_vq(v)s_v\right\rVert^2.
\]
这里正向目标只改变 KL 的积分方向，student 与 score 的定义保持不变。

若数值稳定化把 K3 的函数 $\phi(r)=e^{-r}+r-1$ 改成分段可导的 $\widetilde\phi(r)$，K3$^{++}$ 的数值偏差和梯度偏差分别变成
\[
 \E_p[\widetilde\phi(r_A)]-K,\qquad
 \E_p[(\widetilde\phi(r_A)+\widetilde\phi'(r_A)-r_A)s_A].
\]
公开代码对 log-ratio 与输出均有 clamp，因此无 clipping 恒等式的适用条件是相关 clamp 未激活，或相关偏差项经额外核验为零。

\subsection{数值 estimator 的自动微分梯度}

把已经抽到的动作视为固定索引，对 K1--K4 直接反传，得到
\[
 G_1=s_A,\quad G_2=r_As_A,\quad
 G_3=(1-e^{-r_A})s_A,\quad G_4=r_As_A.
\]
令 $u_q=\E_q[s_A]$，则这些随机梯度相对于 $g$ 的偏差为
\[
 \operatorname{Bias}(G_1)=-g,\quad
 \operatorname{Bias}(G_2)=0,\quad
 \operatorname{Bias}(G_3)=-u_q-g,\quad
 \operatorname{Bias}(G_4)=0.
\]
对应方差逐项为
\begin{align*}
 \Var(G_1)&=\E_p\lVert s_A\rVert^2,\\
 \Var(G_2)&=\E_p[r_A^2\lVert s_A\rVert^2]-\lVert g\rVert^2,\\
 \Var(G_3)&=\E_p[(1-e^{-r_A})^2\lVert s_A\rVert^2]-\lVert u_q\rVert^2,\\
 \Var(G_4)&=\E_p[r_A^2\lVert s_A\rVert^2]-\lVert g\rVert^2.
\end{align*}
例如 $\E_p[G_3]=\E_p[s_A]-\E_p[(q(A)/p(A))s_A]=-u_q$。由此得到两个明确区别：K2 的数值有偏但其固定样本反传梯度无偏；K1 与 K3 的数值无偏但其裸反传梯度通常有偏。这些期望在每次抽样均使用当前 $p$ 时成立；不能将固定样本的导数与包含采样测度导数的 $\nabla_\theta\E_{p_\theta}$ 混用。

确定性 Top-$J$ 数值目标的梯度也可以显式计算。记
\[
 d_P=\nabla_\theta P=\sum_{v\in S}p(v)s_v,
 \qquad \ell_{\tail}=\log(\tau_p/\tau_q).
\]
由乘积法则与 $\nabla_\theta\tau_p=-d_P$，有
\begin{align*}
 \nabla_\theta H&=h_g+d_P,\\
 \nabla_\theta(H+M)&=h_g-\ell_{\tail}d_P=h_g+\tau_pb_{\mathrm{TA}},\\
 \nabla_\theta K_J^{\mathrm{norm}}
 &=\frac{h_g}{P}-\frac{H}{P^2}d_P.
\end{align*}
所以三者相对于完整 KL 梯度的偏差依次为
\[
 d_P-\tau_p\mu_g,\qquad
 \tau_p(b_{\mathrm{TA}}-\mu_g),\qquad
 \frac{h_g}{P}-\frac{H}{P^2}d_P-g,
\]
方差均为零。尤其 $h_g$ 是停止系数的头部 score 梯度，$\nabla H=h_g+d_P$ 是头部数值的真实梯度；两者必须分别命名。

\subsection{停止梯度 surrogate 的显式构造}

为把数学上的梯度向量对应到可反传的标量，记 $\operatorname{sg}$ 为停止梯度算子。固定抽样动作和事件，定义
\begin{align*}
 L_{\mathrm{HETS}}
 &=\sum_{v\in S}\operatorname{sg}(p(v)r_v)\log p(v)
   +J_A\operatorname{sg}(r_A)\log p(A),\\
 L_{\mathrm{Res}}
 &=L_{\mathrm{HETS}}
   +\operatorname{sg}((\tau_p-J_A)\ell_{\tail})\log\tau_p.
\end{align*}
逐项反传，得到
\begin{align*}
 \nabla L_{\mathrm{HETS}}&=h_g+J_Ar_As_A,\\
 \nabla L_{\mathrm{Res}}
 &=h_g+J_Ar_As_A+(\tau_p-J_A)\ell_{\tail}\nabla\log\tau_p\\
 &=h_g+\tau_pb_{\mathrm{TA}}+J_A(r_As_A-b_{\mathrm{TA}}).
\end{align*}
这两条等式依赖停止整个系数的梯度，包括其中的 $p(v)$、$r_v$、$\tau_p$ 和 $\ell_{\tail}$；若遗漏停止操作，就会引入额外导数。

若希望 forward 展示某个数值估计量 $F$，同时 backward 使用 $L$ 的梯度，可定义
\[
 \widetilde L=\operatorname{sg}(F)+L-\operatorname{sg}(L).
\]
逐样本有 $\widetilde L=F$，但 $\nabla\widetilde L=\nabla L$。因此 forward 的数值偏差和 backward 的梯度偏差可以独立设计，审计时必须同时给出 $F$ 和 $L$。

\subsection{K4 尾部反传与 HETS 的精确关系}
\label{sec:kl-k4-hets}

前文 HETS 数值估计量在项目中也称为 EMA，其公式为 $\widehat K^{\mathrm{HETS}}=H+J_Ar_A$。若头部按真实数值 $H$ 求导，而尾部使用本章明确命名的 K4$^{\mathrm{live}}$ 构造，则自动微分得到
\[
 \widehat g^{\mathrm{EMA\text{-}K4}}
 =h_g+d_P+J_A(r_A+1)s_A,\qquad d_P=\nabla_\theta P.
\]
这与 HETS 的 $h_g+J_Ar_As_A$ 相差 $d_P+J_As_A$。因为 $\E_p[J_As_A]=\nabla_\theta\tau_p=-d_P$，差值期望为零，但样本和方差均不同。

\begin{theorem}[头尾 score 常数的匹配]
对固定实数 $\alpha$，定义
\[
 \widehat g_\alpha
 =h_g+\alpha d_P+J_A(r_A+\alpha)s_A.
\]
则它对 $g$ 无偏，且
\begin{equation}
 \Var(\widehat g_\alpha)
 =\sum_{v\notin S}p(v)(r_v+\alpha)^2\lVert s_v\rVert^2
  -\lVert\tau_p\mu_g-\alpha d_P\rVert^2.
 \label{eq:kl-alpha-var}
\end{equation}
$\alpha=0$ 是 HETS，$\alpha=1$ 是真实头部数值加 masked K4$^{\mathrm{live}}$ 尾部的反传。
\end{theorem}

\begin{proof}
期望为
\[
 h_g+\alpha d_P+\tau_p\mu_g+\alpha\sum_{v\notin S}p(v)s_v
 =h_g+\tau_p\mu_g=g.
\]
确定性头部不贡献协方差；随机尾部的均值为 $\tau_p\mu_g-\alpha d_P$，二阶矩为式~\eqref{eq:kl-alpha-var} 的求和项，两者相减得到方差。
\end{proof}

如果误把真实头部梯度 $h_g+d_P$ 与尾部的 $J_Ar_As_A$ 拼接，得到的期望为 $g+d_P$，偏差为 $d_P$。若头部停止系数为 $h_g$，尾部却使用 $J_A(r_A+1)s_A$，偏差为 $-d_P$。两端的 score 常数必须匹配；仅凭“精确头部加 sampled tail”这句话无法判断无偏性。

EMA-PG 固定版本的 \texttt{compute\_memory\_efficient\_kl} 在 \texttt{full\_reverse} 路径中直接反传头部 $H$，尾部使用停止 log-ratio 的 K4。因此在共同固定支持、无 clipping 且尾部测度校正正确时，forward 无偏，backward 的偏差为 $d_P$。actor 传入的 log-importance weight 已 detach，不能通过它的导数消掉 $d_P$。\footnote{同一固定版本的 actor 路径见 \href{https://github.com/LunjunZhang/ema-pg/blob/f9b977632ad948ce935432a517c22291bfaba562/search/verl/workers/actor/dp_actor.py}{\texttt{dp\_actor.py}}。}本章所证明的无偏 EMA 数值结构与论文 HETS 的向量结构，不能据名称相近就认定为该路径的逐样本反传行为。

本仓库的 \path{tests/test_kl_estimators.py} 核验单 token K4$^{\mathrm{live}}$ 的数值与梯度，对应 $(r_A+1)s_A$。\path{experiments/analysis/ema_topk_k_analysis.py} 中的结构统计采用 $h_g+d_P$ 头部与 $r_As_A$ 尾部，其期望具有上述 $d_P$ 偏差；这些统计不能作为正确匹配头尾常数的估计量或 HETS 的无偏性证据。\path{src/mirroropd/training/losses.py} 的 OPD baseline 使用停止梯度的 Tail-aware 标量，属于 scalar baseline 路径。本章的数学分类分别描述这三个实际对象。

\section{采样设计：on-policy、尾部条件采样与旧 rollout}
\label{sec:kl-sampling}

总定理从一开始区分 $\nu$ 与 $\rho$。本节把已有结论用于实际采样设计，所有积分目标和 log-ratio 保持不变。

\subsection{重要性校正与最优尾部常控制}
\begin{corollary}[直接重要性采样]\label{thm:kl-importance}
总定理取 $a=0$，便得到
\[
 \widehat I=\frac{\nu(A)}{\rho(A)}f_A,\qquad
 \E_\rho\widehat I=I_\nu(f),\qquad
 \Var_\rho(\widehat I)
 =\sum_v\frac{\nu(v)^2}{\rho(v)}\|f_v\|^2-\|I_\nu(f)\|^2,
\]
其中求和只作用于 $\rho(v)>0$ 的非零贡献。对反向 KL 的数值和梯度，分别取 $(\nu,f)=(p,r)$、$(p,rs)$；对正向目标，分别取 $(q,-r)$、$(q,-s)$。
\end{corollary}

尾部常控制的近似函数为头部 $a_v=f_v$、尾部 $a_v=b$。本小节比较所有可能的常控制，统一要求 $\rho(v)>0$ 对每个尾部动作成立；精确头部允许 $\rho(v)=0$。这个条件保证二次型对任意 $b$ 都有定义。定义
\[
 h_f=\sum_{v\in S}p(v)f_v,\qquad
 \mu_f=\E_{p_\tail}f_A,\qquad w_v=p(v)/\rho(v).
\]
总定理直接给出
\begin{equation}\label{eq:kl-is-generic}
 \widehat I_\rho(b)=h_f+\tau_pb+J_Aw_A(f_A-b),
\end{equation}
\begin{equation}\label{eq:kl-is-generic-var}
 \Var_\rho(\widehat I_\rho(b))
 =\sum_{v\notin S}\frac{p(v)^2}{\rho(v)}\|f_v-b\|^2
   -\tau_p^2\|\mu_f-b\|^2.
\end{equation}
为了与各估计量的标签对应，标量代入 $f=r,b=\beta$ 的式子记为
\begin{equation}\label{eq:kl-is-scalar}
 \widehat K_\rho(\beta)=H+\tau_p\beta+J_Aw_A(r_A-\beta),
\end{equation}
向量代入 $f=rs$ 的式子记为
\begin{equation}\label{eq:kl-is-vector}
 \widehat g_\rho(b)=h_g+\tau_pb+J_Aw_A(r_As_A-b).
\end{equation}
式~\eqref{eq:kl-is-generic-var} 的两次代入就是它们各自的方差；两种值域无需再证明两遍重要性采样。

\begin{proposition}[任意采样测度下的最优尾部常控制]
记
\[
 W=\sum_{v\notin S}\frac{p(v)^2}{\rho(v)},\qquad
 U_f=\sum_{v\notin S}\frac{p(v)^2}{\rho(v)}f_v,\qquad T_f=\tau_p\mu_f.
\]
若 $W>\tau_p^2$，则
\[
 b_\rho^*=\frac{U_f-\tau_pT_f}{W-\tau_p^2}.
\]
若 $W=\tau_p^2$，尾部控制逐样本抵消，方差与 $b$ 无关。
\end{proposition}
\begin{proof}
展开式~\eqref{eq:kl-is-generic-var}，与 $b$ 有关的项为
\[
 -2\langle U_f-\tau_pT_f,b\rangle+(W-\tau_p^2)\|b\|^2.
\]
当二次项系数严格为正时配方即得最优值。由 Cauchy--Schwarz，
\[
 \tau_p^2
 \le\left(\sum_{v\notin S}\frac{p(v)^2}{\rho(v)}\right)
      \left(\sum_{v\notin S}\rho(v)\right)\le W.
\]
等号要求 $\rho$ 全部质量在尾部，且与 $p_\tail$ 成比例，此时 $\rho=p_\tail$、$J_Aw_A=\tau_p$，式~\eqref{eq:kl-is-generic} 的控制项直接抵消。当 $\rho=p$ 时，$W=\tau_p$、$U_f=T_f$，最优控制退化为 $\mu_f$。
\end{proof}

\subsection{尾部条件采样与真正的分层方差减少}
若能够单独从 $p_\tail$ 采样 $B_1,\ldots,B_n$，则
\[
 \widehat I_\tail=h_f+\frac{\tau_p}{n}\sum_{i=1}^n f_{B_i},
 \qquad
 \Var(\widehat I_\tail)=\frac{\tau_p^2}{n}\Var_{p_\tail}(f).
\]
证明是在总定理中取 $\rho=p_\tail$，再独立平均。常控制 $\beta$ 或 $b$ 逐样本完全抵消，故此时 ResOPD 型的尾部常量控制没有额外作用。相同 $n$ 下，on-policy 头尾控制的方差为
\[
 \frac1n\{\tau_p\Var_{p_\tail}(f)+P\tau_p\|\mu_f-b\|^2\}.
\]
它与条件尾部采样的差为
\[
 \frac{P\tau_p}{n}\{\Var_{p_\tail}(f)+\|\mu_f-b\|^2\}\ge0.
\]
这是数学上的同样本数比较。实际采样器是否能取得尾部条件样本、是否需要额外 teacher 查询和拒绝采样，决定成本比较。

更一般地，若对每个分区单元独立抽样 $n_C$ 次，则
\[
 \widehat I_{\mathrm{strat}}
 =\sum_CP_C\frac1{n_C}\sum_{i=1}^{n_C}f_{B_{C,i}},\qquad
 \Var(\widehat I_{\mathrm{strat}})
 =\sum_C\frac{P_C^2}{n_C}\Var_{p_C}(f).
\]
这类分层积分直接消去了单元之间的抽样质量波动。给定连续预算 $\sum_Cn_C=N$、定义 $\sigma_C^2=\Var_{p_C}(f)$，Cauchy--Schwarz 给出
\[
 \sum_C\frac{P_C^2\sigma_C^2}{n_C}
 \ge\frac{(\sum_CP_C\sigma_C)^2}{N},
 \qquad n_C^*=\frac{NP_C\sigma_C}{\sum_DP_D\sigma_D}.
\]
零方差单元可以精确返回其常值；整数预算需另外取整。这里优化的是协方差迹，向量情形也使用相同 $\sigma_C$。

\subsection{截断与旧 rollout 的偏差由残差决定}
若仍以单位权重使用头尾控制，实际 $A\sim\rho$，总偏差定理给出
\begin{equation}\label{eq:kl-wrong-measure-bias}
 \E_\rho[\widehat I(b)]-I_p(f)
 =\sum_{v\notin S}(\rho(v)-p(v))(f_v-b).
\end{equation}
令 $\tau_\rho=\rho(S^c)$、$\mu_{f,\rho}=\E_\rho[f_A\mid J_A=1]$，条件协方差分解给出
\begin{equation}\label{eq:kl-wrong-measure-var}
 \Var_\rho(\widehat I(b))
 =\tau_\rho\Var_{\rho,\tail}(f_A)
 +(1-\tau_\rho)\tau_\rho\|\mu_{f,\rho}-b\|^2.
\end{equation}
因此更低的实际抽样方差与非零目标偏差可以同时出现。

旧 rollout 应使用实际生成概率 $\rho$ 校正；当前完整 softmax 只定义目标 $p$。温度变换保留支持时，若 $\rho$ 已知便可校正。top-$p$/top-$k$ 删除未知尾部动作时，总定理中的缺失支持项
\[
 -\sum_{v:\rho(v)=0}p(v)(f_v-a_v)
\]
给出无法由有限重要性权重恢复的贡献。只有已精确覆盖的动作残差为零时，支持缺失才不影响该目标。

若把正确权重 $w_v$ 截断为 $\bar w_v$，额外偏差为
\[
 \sum_{v:\rho(v)>0}\rho(v)(\bar w_v-w_v)(f_v-a_v).
\]
这也是同一个偏差定理的特例。通过近似减小未知动作上的残差，是缓解权重误差的方法；它不改变对支持和生成测度的要求。

\subsection{正向 Top-$J$：只替换积分权重与函数}

使用完全相同的 $p,q,S,\tau_p,\tau_q$，定义正向头部
\[
 H^{\mathrm F}=-\sum_{v\in S}q(v)r_v,\qquad
 B^{\mathrm F}=\tau_q\KL(q_{\tail}\Vert p_{\tail}).
\]
逐项将尾部 log-ratio 拆成质量与条件形状，得到
\[
 K^{\mathrm F}=H^{\mathrm F}-\tau_q\ell_{\tail}+B^{\mathrm F}.
\]
因此正向 Truncated、Tail-aware、Normalized 三个数值目标分别为
\[
 H^{\mathrm F},\qquad
 H^{\mathrm F}-\tau_q\ell_{\tail},\qquad
 \frac{H^{\mathrm F}}{Q}+\log\frac{P}{Q}.
\]
相对完整正向 $K^{\mathrm F}$ 的偏差分别为
\[
 \tau_q\ell_{\tail}-B^{\mathrm F},\qquad
 -B^{\mathrm F},\qquad
 \frac{H^{\mathrm F}}{Q}+\log\frac{P}{Q}-K^{\mathrm F},
\]
且固定 prefix 下方差均为零。这里 $B^{\mathrm F}$ 与前文 $B$ 分别对应不同方向的尾部条件 KL，上标始终保留。

对标量尾部控制 $\beta$，完整正向 KL 的无偏校正为
\[
 \widehat K^{\mathrm F}_\rho(\beta)
 =H^{\mathrm F}+\tau_q\beta
   +J_A\frac{q(A)}{\rho(A)}(-r_A-\beta).
\]
令 $\mu_r^{\mathrm F}=\E_{q_{\tail}}[-r_A]$，则其方差为
\[
 \Var_\rho(\widehat K^{\mathrm F}_\rho(\beta))
 =\sum_{v\notin S}\frac{q(v)^2}{\rho(v)}(-r_v-\beta)^2
   -\tau_q^2(\mu_r^{\mathrm F}-\beta)^2.
\]
证明是定理~\ref{thm:kl-importance} 在尾部对积分分布 $q$ 的直接应用：修正项均值是 $\tau_q(\mu_r^{\mathrm F}-\beta)$，恰好补齐尾部。

梯度方面，定义 $h_g^{\mathrm F}=-\sum_{v\in S}q(v)s_v$。因为 $q$ 固定，逐项求导给出
\begin{align*}
 \nabla H^{\mathrm F}&=h_g^{\mathrm F},\\
 \nabla(H^{\mathrm F}-\tau_q\ell_{\tail})
 &=h_g^{\mathrm F}-\tau_q\bar s_{\tail},\\
 \nabla\left(\frac{H^{\mathrm F}}{Q}+\log\frac P Q\right)
 &=\frac{h_g^{\mathrm F}}{Q}+\frac{d_P}{P}.
\end{align*}
三者的梯度偏差分别是上述向量减去 $g^{\mathrm F}$，梯度方差均为零。这里 $\bar s_{\tail}=\nabla\log\tau_p$ 是 student 尾部平均 score；它不能替换为未知的 teacher 尾部平均 score。特别地，正向 Tail-aware 梯度偏差为
\[
 \sum_{v\notin S}q(v)s_v-\tau_q\bar s_{\tail}.
\]
正向头部梯度与反向头部梯度的差异来自积分权重：$H^{\mathrm F}$ 的权重 $q$ 冻结，因此没有反向 $\nabla H$ 中的 $d_P$ 项。

正向无偏向量控制同样使用原来的向量符号 $b$：
\[
 \widehat g^{\mathrm F}_\rho(b)
 =h_g^{\mathrm F}+\tau_qb
   +J_A\frac{q(A)}{\rho(A)}(-s_A-b).
\]
令 $\mu_g^{\mathrm F}=\E_{q_{\tail}}[-s_A]$，有
\[
 \Var_\rho(\widehat g^{\mathrm F}_\rho(b))
 =\sum_{v\notin S}\frac{q(v)^2}{\rho(v)}\lVert-s_v-b\rVert^2
  -\tau_q^2\lVert\mu_g^{\mathrm F}-b\rVert^2.
\]
所有正向估计的支持覆盖要求作用于未知的 $q$ 尾部。这样，BPO 的正向 student--anchor center 与 OPD 的反向 student--teacher KL 都使用同一套采样与误差定义。

\section{OPD 使用方式与误差传递}
\label{sec:kl-opd-roles}
\subsection{Score baseline 是积分为零的控制}

本小节默认 $\rho=p$。vOPD 在总定理中对应 $f_v=r_vs_v$、$a_v=c(x)s_v$。由于 $I_p(a)=c(x)I_p(s)=0$，它直接成为 $(r_A-c)s_A$。对其他覆盖残差支持的采样分布，正确形式是 $(p(A)/\rho(A))(r_A-c)s_A$。因此 baseline 是否对 KL 数值无偏，不是这个构造的无偏条件；控制积分为零才是条件。下面求出的最优 $c$ 就是 Gram 定理的单列控制解。


Sampled-token estimator 为
\[
 \widehat g^{\mathrm{sample}}=r_As_A.
\]
它无偏，且
\begin{equation}
 \Var(\widehat g^{\mathrm{sample}})
 =\E_p[r_A^2\lVert s_A\rVert_2^2]-\lVert g\rVert_2^2.
 \label{eq:chapter3-sample-grad-var}
\end{equation}

对任意在抽样前固定的 scalar $c(x)$，vOPD 型 estimator 为
\begin{equation}
 \widehat g^{\mathrm{vOPD}}_c=(r_A-c)s_A.
 \label{eq:chapter3-vopd}
\end{equation}

\begin{proposition}[固定 scalar baseline 的无偏性与最优值]
\label{prop:chapter3-vopd}
对任意 action-independent $c(x)$，$\widehat g^{\mathrm{vOPD}}_c$ 无偏。其二阶矩为
\[
 \E\lVert\widehat g^{\mathrm{vOPD}}_c\rVert_2^2
 =\E_p[(r_A-c)^2\lVert s_A\rVert_2^2],
\]
当 $\E_p\lVert s_A\rVert^2>0$ 时，使该二阶矩最小的 scalar 是
\begin{equation}
 c^*=\frac{\E_p[r_A\lVert s_A\rVert_2^2]}{\E_p[\lVert s_A\rVert_2^2]}.
 \label{eq:chapter3-optimal-score-baseline}
\end{equation}
一般 $c^*\neq K$，因此用 KL 数值 $K$ 作为 vOPD baseline 是合法的中心化，但不必是 score-gradient 方差最优的 baseline。
若 $\E_p\lVert s_A\rVert^2=0$，score 在 $p$ 下几乎处处为零，所有 baseline 对应的梯度均为零。
\end{proposition}

\begin{proof}
\[
 \E_p[(r_A-c)s_A]=\E_p[r_As_A]-c\E_p[s_A]=g.
\]
对二阶矩关于 $c$ 求导并令其为零，得到式~\eqref{eq:chapter3-optimal-score-baseline}。
\end{proof}

进一步记 $D_s=\E_p\lVert s_A\rVert^2>0$，展开平方并配方，得到精确的方差表达式
\begin{align*}
 \Var(\widehat g_c^{\mathrm{vOPD}})
 &=\E_p[r_A^2\lVert s_A\rVert^2]
   -2c\E_p[r_A\lVert s_A\rVert^2]+c^2D_s-\lVert g\rVert^2\\
 &=\Var(\widehat g_{c^*}^{\mathrm{vOPD}})+D_s(c-c^*)^2.
\end{align*}
这给出任意固定 KL proxy 作为 baseline 的代价，而无需假定该 proxy 对 $K$ 无偏。

\begin{proposition}[独立随机 KL baseline 的误差传递]
令 $C$ 与当前动作 $A$ 条件独立，$m=\E[C]$、$\sigma_C^2=\Var(C)$，且反传时停止 $C$ 的梯度。则 $(r_A-C)s_A$ 对 $g$ 无偏，且
\[
 \Var((r_A-C)s_A)
 =\Var((r_A-m)s_A)+D_s\sigma_C^2.
\]
若写 $m=K+\delta$，则
\[
 \Var((r_A-C)s_A)
 =\Var(\widehat g_{c^*}^{\mathrm{vOPD}})
  +D_s\{(K+\delta-c^*)^2+\sigma_C^2\}.
\]
\end{proposition}

\begin{proof}
条件于 $A$，$\E_C[(r_A-C)^2]=(r_A-m)^2+\sigma_C^2$。乘以 $\lVert s_A\rVert^2$ 后取 $A$ 的期望，并减去相同的 $\lVert g\rVert^2$，即可得到结论。
\end{proof}

所以在 vOPD 中只改进后面的 KL 数值估计时，应比较 scalar baseline 的均值与方差；当两个 baseline 均对 $K$ 无偏且独立时，梯度方差之差恰好是 $D_s$ 乘以 scalar 方差之差。ResOPD 的向量控制方差结论不能直接用于这个比较。

如果 baseline 是与当前 action 共同随机化的 $C_A$，则上述证明失效：偏差为
\begin{equation}
 \operatorname{Bias}(\widehat g_C)=-\E[C_As_A].
 \label{eq:chapter3-random-baseline-bias}
\end{equation}
因此随机 baseline 保持无偏的充要条件是 $\E[C_As_A]=0$。独立样本或 leave-one-out/cross-fitting 是保证这个条件的构造方式；仅有 $\E[C_A]=K$ 不够。

同样本 K1 baseline 提供一个极端反例：取 $C_A=r_A$，则 $(r_A-C_A)s_A=0$，偏差恰好为 $-g$，尽管 $\E[C_A]=K$。数值无偏不能消除 baseline 与 score 的相关性。

\subsection{旧 rollout 下的最优 baseline}

若 $A\sim\rho$，且 $\rho$ 覆盖 $p$ 的支持，令 $w_A=p(A)/\rho(A)$。控制 $c s_A$ 的正确修正仍然无偏，因为
\[
 \E_\rho[w_A(r_A-c)s_A]=I_p(rs)-cI_p(s)=g.
\]
但是其二阶矩使用 $p^2/\rho$ 权重。定义
\[
 D_{s,\rho}=\sum_v\frac{p(v)^2}{\rho(v)}\|s_v\|^2.
\]
若该值为正，对二阶矩关于 $c$ 配方得到
\[
 c_\rho^*=
 \frac{\sum_v p(v)^2\rho(v)^{-1}r_v\|s_v\|^2}{D_{s,\rho}},
\]
\[
 \Var_\rho(w_A(r_A-c)s_A)
 =\Var_\rho(w_A(r_A-c_\rho^*)s_A)
  +D_{s,\rho}(c-c_\rho^*)^2.
\]
这直接来自 Gram 定理取控制列 $w_As_A$。若 $D_{s,\rho}=0$，所有 score 均为零，baseline 不影响结果。

对与当前动作独立的随机 baseline $C$，令 $m=\E C$、$\sigma_C^2=\Var C$。条件于动作展开平方，得到
\[
 \Var_\rho(w_A(r_A-C)s_A)
 =\Var_\rho(w_A(r_A-m)s_A)+D_{s,\rho}\sigma_C^2.
\]
因此两个独立、同均值的 KL baseline 的比较，只需比较其标量方差，梯度方差差为 $D_{s,\rho}$ 乘以该差；当前 on-policy 情形是 $\rho=p$ 的特例。若 baseline 与动作相关，无偏性的条件也相应变成 $\E_\rho[w_AC_As_A]=0$。

\subsection{Sampled OPD 与 vOPD}

固定 teacher 作为 $q$，sampled OPD 取停止系数 $r_A$，其最小化梯度为 $r_As_A$；若使用 reward 最大化记号，reward 是 $-r_A$，两种表述符号相反而更新方向相同。vOPD 取 $(r_A-c(x))s_A$，其中 $c(x)$ 可取完整 $K$、Tail-aware 下界或独立随机 KL baseline。第三章已经证明：action-independent baseline 的数值偏差不产生梯度偏差；同 action 的无偏 KL 估计反而可以产生梯度偏差。

本仓库生产 OPD baseline 对 $q=\pi_T$ 使用 $H+\tau_p\log(\tau_p/\tau_q)$，并在梯度系数中停止。它是完整 KL 的有偏数值下界，但作为 action-independent baseline，在正确采样测度或正确重要性校正下保持 OPD 梯度无偏。若 rollout 来自冻结 $\pi_k$，则此处 $\rho=\pi_k$，需要 $p/\rho$；baseline 本身使用 teacher 作为 $q$，二者不可混用。

\subsection{Mirror 目标中的 student--anchor KL}

对 MDPO，令本节参考 $q=\pi_k$，固定 reward 为前两章的 $r_k(x,v)$。负 mirror 目标为
\[
 \mathcal L_{\mathrm{MDPO}}
 =\eta^{-1}K-\E_p[r_k(x,A)].
\]
逐项求导并利用 score 均值为零，得到
\[
 \nabla\mathcal L_{\mathrm{MDPO}}
 =\E_p[u_A^{\mathrm{MDPO}}s_A],\qquad
 u_v^{\mathrm{MDPO}}=\eta^{-1}r_v-r_k(x,v).
\]
在冻结 rollout $\rho=\pi_k$ 下，无偏估计为
\[
 \widehat g_{\mathrm{MDPO}}=\frac{p(A)}{\rho(A)}u_A^{\mathrm{MDPO}}s_A,\qquad
 \Var_\rho(\widehat g_{\mathrm{MDPO}})
 =\sum_v\frac{p(v)^2}{\rho(v)}(u_v^{\mathrm{MDPO}})^2\lVert s_v\rVert^2
  -\lVert\nabla\mathcal L_{\mathrm{MDPO}}\rVert^2.
\]
如果 $\lVert u_v^{\mathrm{MDPO}}s_v\rVert\leq C$，则二阶矩满足
\begin{equation}
 \E_\rho\lVert\widehat g_{\mathrm{MDPO}}\rVert^2
 \leq C^2\sum_v\frac{p(v)^2}{\rho(v)}
 =C^2\{1+\chi^2(p\Vert\rho)\}.
 \label{eq:mdpo-variance}
\end{equation}
这说明 stale rollout 与当前 student 分离时的重要性权重代价。

BPO 的 student center 是相反方向的 $\KL(\pi_k\Vert p)$，在本节 $q=\rho=\pi_k$ 下就是 $K^{\mathrm F}$；teacher center 是冻结的 $\KL(\pi_k\Vert\pi_T)$。PMD-Mean 的目标使用冻结 reward 中心而不含同样的 student KL center。它们的原目标就在 anchor 测度下定义，不能因 MDPO 需要 $p/\rho$，便给 Mean 或 BPO 机械乘上该比率。center 放入平方损失后还会出现同样本相关项，须使用第二章的有限样本与留一修正；一个 KL scalar estimator 的无偏性不替代整个目标梯度的证明。

\subsection{确定性代理的平方误差}

记 $M=\tau_p\ell_{\tail}$、$B=\tau_p\KL(p_{\tail}\Vert q_{\tail})$。Tail-aware 与截断的平方误差差为
\[
 (K-H)^2-\{K-(H+M)\}^2=(M+B)^2-B^2=M(M+2B).
\]
因此 $M\geq0$ 时 Tail-aware 的标量平方误差不大于截断；$M<0$ 时删除的质量项与形状项可抵消，零方差确定性 proxy 的排序不能只按是否保留更多项判断。

支持细分的 KL 单调性已经由第~\ref{sec:kl-partition} 节的一般分区定理证明；该性质控制确定性粗化偏差，不能推出 HETS 或 ResOPD 的抽样方差单调下降。

\subsection{Normalized 与 Truncated 的 logit 行为}

令 $p=\operatorname{softmax}(\theta)$，其 score 为 $s_{v,j}=\mathbf1\{v=j\}-p(j)$。Normalized Top-$J$ 只依赖头部相对 logits，故对 $j\notin S$ 有
\[
 \frac{\partial K^{\mathrm{norm}}_J}{\partial\theta_j}=0.
\]
截断数值的真实梯度则满足
\[
 \frac{\partial H}{\partial\theta_j}=-p(j)(H+P),\qquad j\notin S.
\]
这由 $\nabla H=\sum_{v\in S}p(v)(r_v+1)s_v$ 逐坐标代入得到。当 $H+P>0$ 时，梯度下降增加尾部 logits。停止系数的头部梯度 $h_g$ 省去了 $P$ 对应的边界项，但仍缺失完整尾部。因此这两种截断训练路径也必须分别报告。

\subsection{固定 prefix 误差与访问测度下误差}

令 $x\sim d$，条件偏差为 $\beta(x)=\E[\widehat K\mid x]-K(x)$，条件方差为 $V(x)$。则
\[
 \E_{x,A}[\widehat K-K]=\E_d[\beta(x)],\qquad
 \E_{x,A}[(\widehat K-K)^2]=\E_d[\beta(x)^2+V(x)].
\]
而原始 estimator 的总体方差是
\[
 \Var_{x,A}(\widehat K)
 =\E_d[V(x)]+\Var_d(K(x)+\beta(x)).
\]
这三个统计量分别回答总体有符号误差、逐 prefix 均方误差和 estimator 自身的总体波动。确定性 Tail-aware 有 $V(x)=0$，但 $\operatorname{RMSE}=\sqrt{\E_d[B(x)^2]}$，一般不等于 $|\E_d[-B(x)]|$。向量估计量将平方替换为范数平方、方差替换为协方差迹，结论保持相同形式。访问测度 $d$ 在不同 outer 轮次的变化由第四章单独处理。

\section{ResOPD 实现与论文实验}
\label{sec:resopd-implementation}

本节沿用前文的 $r_v$、score $s_v$、头部向量 $h_g$、尾部均值 $\mu_g$ 与事件 $J_A$，将控制族对应到公开训练实现。\footnote{论文来源为 \href{https://arxiv.org/abs/2610.04882}{ResOPD}；本仓库保存的全文为 \path{docs/papers/2610.04882_zh_CN/}。代码来源为 \href{https://github.com/InternLM/ResOPD}{InternLM/ResOPD}。}

\subsection{代码调用链与停止梯度}

开源仓库没有单独的 ``HETS'' loss mode；HETS 是论文中的比较基线。发布代码注册的唯一新模式是 ``resopd''，而 sampled-token、Top-$k$ 和普通 ``k1/k2/k3'' 路径属于其他 loss mode。

teacher 路径从 \path{verl/experimental/teacher_loop/teacher_manager.py} 开始。选择 ``loss\_mode=resopd'' 后，vLLM 必须返回 teacher Top-$J$ 的 token ids/log-probs，以及实际 student token 的 teacher log-prob。

\path{verl/workers/rollout/vllm_rollout/prompt_logprobs.py} 保存两组字段：
\begin{itemize}
\item 头部支持：\texttt{prompt\_ids} 与 \texttt{prompt\_logprobs}。
\item 实际动作：\texttt{sampled\_prompt\_ids} 与 \texttt{sampled\_prompt\_logprobs}。
\end{itemize}
实际动作的 teacher log-prob 即使落在 Top-$J$ 之外也必须保留。

student logits 经 \path{verl/trainer/distillation/fsdp/losses.py} 对齐后，\path{verl/trainer/distillation/losses.py} 将四个 payload 送入 \path{verl/trainer/distillation/resopd.py} 的 \texttt{resopd\_loss}。该 kernel 的数学对应关系如下：
\begin{enumerate}
\item 完整 student logits 做 log-softmax，得到头部 $p(v)$、$r_v$ 和 $P$；teacher Top-$J$ log-prob 在进入计算时 detach。
\item 从实际 token 得到 $r_A$，并用 \texttt{student\_tail\_mass}、\texttt{teacher\_tail\_mass} 构造 $\log(\tau_p/\tau_q)$。
\item 用 \texttt{support\_score} 计算 $b_{\mathrm{TA}}$，并令 $\texttt{centered\_event}=\mathbf 1\{A\in S\}-P$。
\item 通过 \texttt{surrogate - surrogate.detach()} 保留 forward 数值、替换 backward，得到式~\eqref{eq:chapter3-control-family} 中 $b=b_{\mathrm{TA}}$ 的特例。
\end{enumerate}
实现中的 teacher 数值、事件系数和控制项均停止梯度。仓库 docstring 明确说明返回值不是确定性 KL 目标，单个 forward value 可以为负；应审计 backward 梯度、期望和协方差，而不是把 loss 平均值解释成完整 KL。

无偏性要求 $A$ 来自当前完整 student 分布，或按本章给出的重要性权重校正。发布脚本使用温度 $1$、\texttt{top\_p=1}、\texttt{top\_k=-1}。论文的主要数学实验使用温度 $1$、\texttt{top\_p=0.95}、\texttt{top\_k=20}；全支持审计与跨领域实验采用完整祖先采样。这些设置必须分别报告：重新计算未截断 logits 能恢复目标概率 $p$，不能把已生成动作的测度 $\rho$ 改回 $p$。

\subsection{公开实验与 baseline}

论文的基线可按目标分成三组。第一组是有偏确定性 Top-$k$ 方法：Teacher Top-$k$ fKL、Teacher Top-$k$ rKL 和 TA-OPD（尾部伪 token 的粗粒度 KL）。第二组是无偏 reverse-KL 梯度 estimator：Sampled-token OPD、教师 Top-$k$ 支持上的 vOPD、HETS 和 ResOPD。第三组是机制消融：支持大小 $k\in\{1,4,16,32,64\}$、完整祖先采样与 top-$p$/top-$k$ 截断、学习率敏感性，以及统一双控制族中的 SC-TA（$\gamma=1$；ResOPD 为 $\gamma=0$）。

冻结 prefix 审计在 GSM8K、HMMT26 和 DeepMath 上枚举完整动作，报告梯度协方差迹和 NSR。Top-$4$ 下，论文报告 ResOPD 相对 sampled-token OPD 降低 $52.5\%$--$75.0\%$ 的协方差迹，并在五个面板中优于 HETS。端到端数学实验使用 Qwen3.5-27B 教师和 Qwen3.5-4B 学生，评估 DeepMath、AIME26、HMMT26、OlymMATH；科学/代码迁移使用 Qwen3-30B-A3B-Instruct 到 Qwen3-4B，评估 GPQA-D、SuperGPQA、LiveCodeBench v5/v6。论文报告数学平均分 ResOPD $50.12\%$、HETS $45.72\%$，迁移平均分 ResOPD $44.15\%$、HETS $43.63\%$。

仓库公开了可直接运行的 DeepMath 子集：1,024 条训练样本和 252 条 held-out 样本；入口为 \path{examples/resopd/run_resopd.sh}，默认 FSDP2 student、vLLM teacher、Top-$4$、full-support rollout。\path{tests/resopd} 包含 estimator、teacher payload 和 verl integration 检查，发布验证记录 22 个 standalone 与 13 个 integration checks。

需要避免的概念混淆是：HETS 是稀疏 rKL 梯度基线，不是 PPO 的 value baseline；ResOPD 的 $b_{\mathrm{TA}}$ 是 control variate。仓库中普通 verl 的 ``k1/k2/k3'' KL penalty、GRPO、RLOO、ReMax 等属于框架通用路径，并非论文中同一 teacher payload 下的 HETS/ResOPD 方差基线。

